\documentclass[11pt,a4paper]{article}
\usepackage[T1]{fontenc}
\usepackage[utf8]{inputenc}
\usepackage{lmodern}
\usepackage{microtype}
\usepackage[margin=27mm,headheight=15pt]{geometry}
\usepackage{amsmath,amssymb,amsthm,mathtools}
\usepackage{booktabs,array,longtable,enumitem}
\usepackage{xcolor}
\usepackage{fancyhdr}
\usepackage{hyperref}
\usepackage{bookmark}
\usepackage{xurl}
\usepackage{listings}% Part II (batch 44)
\definecolor{inkblue}{HTML}{183D59}
\definecolor{linkblue}{HTML}{215F87}
\hypersetup{colorlinks=true,linkcolor=linkblue,citecolor=linkblue,urlcolor=linkblue,
 pdftitle={All-Degree Rigidity of a Weighted Keller Class},
 pdfauthor={ChatGPT; prepared for Vladimir Reshetnikov},
 pdfsubject={Classification, degenerations, and an obstructed double point}}
\pagestyle{fancy}
\fancyhf{}
\fancyhead[L]{\small\color{inkblue}All-degree rigidity of a weighted Keller class}
\fancyhead[R]{\small Research draft}
\fancyfoot[C]{\thepage}
\renewcommand{\headrulewidth}{0.3pt}
\setlength{\parskip}{0.35em}
\setlength{\parindent}{1.2em}
\setlength{\emergencystretch}{2em}
\setlist{itemsep=0.25em,topsep=0.4em}
\numberwithin{equation}{section}
\newtheorem{theorem}{Theorem}[section]
\newtheorem{lemma}[theorem]{Lemma}
\newtheorem{proposition}[theorem]{Proposition}
\newtheorem{corollary}[theorem]{Corollary}
\theoremstyle{definition}
\newtheorem{definition}[theorem]{Definition}
\newtheorem{example}[theorem]{Example}
\theoremstyle{remark}
\newtheorem{remark}[theorem]{Remark}
\newcommand{\Q}{\mathbb{Q}}
\newcommand{\A}{\mathbb{A}}
\newcommand{\Gm}{\mathbb{G}_{\mathrm m}}
\newcommand{\J}{\mathcal{J}}
\newcommand{\Spec}{\operatorname{Spec}}
\newcommand{\diag}{\operatorname{diag}}
\newcommand{\wt}{\operatorname{wt}}
\newcommand{\coeff}[2]{[#1]#2}
\newcommand{\cstar}{c_*}
\newcommand{\eps}{\varepsilon}
\newcommand{\repo}{\texttt{ProveIt}}
\newcommand{\dd}{\mathrm d}
\newcommand{\wtag}{\textup{[write]}}% text added when the report was written into ProveIt
\newcommand{\rpath}[1]{\nolinkurl{#1}}
% Part II (batch 44, manuscript 04)
\newcommand{\ord}{\operatorname{ord}}
\newcommand{\Frac}{\operatorname{Frac}}
\newcommand{\Dcap}{\mathsf D}
\lstset{basicstyle=\small\ttfamily,breaklines=true,columns=fullflexible,
 keepspaces=true,frame=single,rulecolor=\color{black!20},xleftmargin=5pt,xrightmargin=5pt}

\begin{document}
\hypersetup{pageanchor=false}
\begin{titlepage}
\thispagestyle{empty}
\noindent{\color{inkblue}\large RESEARCH REPORT}\par
\vspace{1.7cm}
{\color{inkblue}\Huge\bfseries All-Degree Rigidity\par
\vspace{0.2cm}of a Weighted Keller Class\par}
\vspace{0.65cm}
{\Large Classification, degenerations,\par and an obstructed double point\par}
\vspace{1.1cm}
{\large Prepared for Vladimir Reshetnikov}\par
\medskip
September 24, 2026\par
\vspace{0.9cm}
\begin{abstract}
We classify, without any bound on ordinary degree, a class of polynomial
Keller maps in three variables naturally singled out by the Jacobian project
in the \repo\ repository. The source weights are $(-1,1,2)$ and the target
weights are $(2,1,-1)$. In invariant coordinates $t=xy$ and $v=x^2z$, we
assume that the first two numerators are affine in $v$ and that the third
numerator is affine in $(t,v)$. Every normalized map is either an explicitly
invertible tame map or a member of a two-parameter diagonal orbit of the
known Alp\"oge map. In the latter case its degree profile is forced to be
$(7,6,4)$ and its expanded support has exactly sixteen monomials.

The proof replaces a degree-bounded coefficient search by one-variable
polynomial rigidity. An equivariant source shear extends the classification
 to arbitrary polynomial dependence of $r$ on $t$ when its $v$-coefficient
 is constant, and gives the exact noninvertible degree spectrum
 $\{7,8,10,12,\ldots\}$. We also identify the reduced closure of the nonconstant
branch as an affine plane whose coordinate axes are precisely its
invertible members. Finally, for the normalized degree-seven coefficient
problem with third numerator $1+t+v$, we prove that the entire coefficient
scheme is a double point, $\Spec\Q[\eps]/(\eps^2)$. An explicit nonzero
infinitesimal deformation is obstructed at second order. All finite
identities and obstruction certificates are independently reproduced by an
included exact-arithmetic script.
\end{abstract}
\vfill
\noindent\small
\textbf{Status.} Human-readable proofs and reproducible symbolic certificates;
not a Lean/Rocq formalization or a peer-reviewed publication. The results
extend the inspected repository statements, but publication priority is not
established. This report does not claim a new counterexample to the Jacobian
conjecture or a classification of unrestricted Keller maps.

\smallskip
\noindent\wtag{} Batch~36, manuscript~01, placed in \texttt{1a1396d4d}: a
report of ProveIt's research-report collection that continues the formal
project \rpath{Algebra/JacobianConjecture}. AI-assisted, unrefereed, not
formalized; labels carry the prefix \texttt{wkr:}. The weighted
determinant identity and the collision mechanism are the project's own
(Section~\ref{wkr:sub:project}); Gao's three-dimensional map has
\emph{geometric} degree four (Section~\ref{wkr:sub:gao}). Provenance:
Appendix~\ref{wkr:app:report}. [Added 29 September 2026, batch~44:]
Part~II (Section~\ref{wkr:as:sec:provenance}, labels \texttt{wkr:as:},
not formalized) answers Section~\ref{wkr:as:q:slope}.
\end{titlepage}
\hypersetup{pageanchor=true}

\begingroup
\footnotesize
\setlength{\parskip}{0pt}
\tableofcontents
\endgroup
\newpage

\section{Research question, provenance, and contribution}
\label{wkr:sec:intro}

A useful research question should distinguish a structural theorem from a
large but finite search. The Jacobian component of \repo\ provides a
particularly concrete opportunity: it gives explicit maps, exact collision
certificates, a weighted determinant reduction, and degree-bounded searches
for simpler representatives. We ask whether one of its natural coefficient
classes is rigid \emph{without imposing a degree bound}. We then ask what
remains of the apparent uniqueness when coefficients are allowed to contain
nilpotents. These are separate questions: a unique field-valued solution
need not be a reduced point of its coefficient scheme.

\subsection{The repository baseline}

The inspected snapshot is commit
\begin{center}
\texttt{e21766d04c2b8a9b2cdba0cd43e563024b3bd1b9}.
\end{center}
The relevant files are the Jacobian project README, its research README, and
\texttt{verify\_equivariant\_sparsity.py}; immutable references are supplied
in the bibliography \cite{repo-main,repo-research,repo-sparsity}.

The repository's three-variable map can be written
\begin{equation}\label{wkr:eq:original}
\begin{gathered}
u=1+xy,\qquad h=u^2z+y^2(1+3u),\\
F_0(x,y,z)=\bigl(uh,\ y+3xh,\ x(5-3u-x^2z)\bigr).
\end{gathered}
\end{equation}
Its Jacobian determinant is $-2$, and
\begin{equation}\label{wkr:eq:oldcollision}
F_0(-1,1,5)=F_0(0,-2,-16)=(0,-2,0).
\end{equation}
These identities, their formal verification, and their provenance are
already part of the repository \cite{repo-main}. We verify the displayed
identities again in the companion script but do not present them as new.

The research files restrict a sparsity calculation to maps of ordinary
degree at most seven whose invariant numerators $p,q$ are affine in $v$ and
whose third numerator $r$ is affine in $(t,v)$. After normalizing
$r=1+t+v$, the script proves that eleven optional coefficients cannot vanish,
and identifies the branch $q_{40}=0$ with the normalized map. Its certificate
contains a square $(q_{21}-4)^2$, but the repository discussion does not
identify the entire coefficient scheme as the dual numbers
\cite{repo-research,repo-sparsity}.

\subsection{What this article establishes}

The main result, Theorem~\ref{wkr:thm:classification}, removes the ordinary
degree restriction from that structural class. It also handles constant
$r$ and proves that the two one-term nonconstant cases are impossible. The
proof forces the normalized nonconstant map outright, so it proves
$q_{40}=0$ rather than imposing it.

Theorem~\ref{wkr:thm:triangular} extends the result beyond affine $r$ by an
explicit triangular normal form. Its degree-spectrum corollary determines
all possible maximum degrees in that larger class.
Theorem~\ref{wkr:thm:plane} describes the reduced two-parameter closure and its
invertible boundary. Theorem~\ref{wkr:thm:double} strengthens uniqueness in the
fixed normalized degree-seven slice to an exact scheme presentation, and Theorem~\ref{wkr:thm:torusdouble} identifies the full
scheme on the nonconstant-parameter open locus. These results connect
three levels of information:
\begin{center}
\begin{tabular}{@{}p{0.25\textwidth}p{0.66\textwidth}@{}}
\toprule
Level & Conclusion proved here\\
\midrule
Field-valued maps & An all-degree automorphism/counterexample dichotomy.\\
Reduced families & No deformation after fixing $r=1+t+v$; a parameter plane
before fixing its two nonconstant coefficients.\\
Nonreduced families & One square-zero parameter in the degree-seven slice,
with a nonzero second-order obstruction.\\
\bottomrule
\end{tabular}
\end{center}

\subsection{Relation to contemporary literature and limits}

Shaska studies graded Keller maps, their quotient geometry, and the role of
the sign pattern of weights \cite{shaska}. Gao gives a tangent-sweep
construction and further explicit examples, including a three-dimensional
map of degree four \cite{gao}. Our subject is narrower: rigidity within a
specified invariant-numerator class, not the construction of the known map
or a minimum-degree theorem for all three-variable maps. In particular, our
forced degree seven is compatible with examples outside this class of lower
degree.

The searched sources and the inspected repository justify the comparison
above, not a universal priority claim. The mathematical contribution of this
report is its stated and proved classification and deformation results.
Whether equivalent formulations already occur elsewhere requires further
bibliographic review. Nothing in the proofs depends on accepting a new
external claim about the general Jacobian problem: the particular map,
its determinant, and its collisions are explicit.

\subsection{What the project already contains}\label{wkr:sub:project}
\wtag{} This subsection and the next were added when the manuscript was
written into ProveIt. Two ingredients that the manuscript presents in its
own words are already recorded in the project's research notes
\cite{repo-research}, \rpath{Algebra/JacobianConjecture/Research/README.md},
lines 188--224, at the pin and unchanged at the placement commit:
\begin{itemize}
\item the lift $F=(p(t,v)/x^2,\,q(t,v)/x,\,x\,r(t,v))$ with $t=xy$,
$v=x^2z$, the polynomiality conditions of \eqref{wkr:eq:polynomiality}
in their general form, \emph{the same determinant identity}
\eqref{wkr:eq:J} with the same chain-rule proof in the localization at
$x$ (Proposition~\ref{wkr:prop:jacobian}), and the normalization
\eqref{wkr:eq:normalization} with determinant $-1$ (lines 188--208);
\item the collision mechanism: with $\alpha=[t^2]p$ (the notes write $a$,
which is a family parameter here), $F(0,y,z)=(\alpha y^2+z,y,0)$ and
$F(1,t_0,v_0)=F(0,q(t_0,v_0),p(t_0,v_0)-\alpha q(t_0,v_0)^2)$ whenever
$r(t_0,v_0)=0$ (lines 210--224). The collision \eqref{wkr:eq:newcollision}
is this mechanism at $t_0=0$, $v_0=-1/\gamma$.
\end{itemize}
The manuscript cites the research notes for its baseline but not at these
two points; they are credited here. What is new relative to the project is
the all-degree classification and everything after it: the project's
sparsity certificate (\rpath{verify_equivariant_sparsity.py}, described in
lines 240--248 of the notes) covers ordinary degree at most seven with $q_{40}=0$ left as
the ``remaining branch'', and says it is ``deliberately not advertised as a
global sixteen-monomial theorem''. Corollary~\ref{wkr:cor:support} proves
$q_{40}=0$ and the support profile in all degrees, inside the class only.
For constant $r$ the notes invoke Moh's plane theorem through degree~100
(lines 232--238), in a class that allows $v^2$ terms; case~\ref{wkr:case:tame}
proves tameness in all degrees without Moh, but only for $p,q$ affine in
$v$, so the two statements sit side by side.

\emph{Formal status.} The project's Lean and Rocq/Coq developments
formalize the map $F_0$ (Lean \texttt{counterexample},
\rpath{Algebra/JacobianConjecture/Lean/JacobianConjecture/Counterexample.lean},
lines 95--101), its determinant $-2$ (\texttt{jacobianDet\_counterexample}),
the collision \eqref{wkr:eq:oldcollision} (\texttt{collision}), the
weight-$(-1,1,2)$ torus action (\texttt{counterexample\_scaling},
\rpath{Algebra/JacobianConjecture/Lean/JacobianConjecture/Scaling.lean}), the mirror symmetry and the refutation of the
three-dimensional conjecture; the Coq development proves the
corresponding statements. None of the weighted determinant identity,
the normalization \eqref{wkr:eq:originalnormalization}, the family
$F_{a,b}$ or any theorem of this report is formalized.

\subsection{Gao's three-dimensional map has geometric degree four}\label{wkr:sub:gao}
\wtag{} The ``three-dimensional map of degree four'' of \cite{gao} cited
above and in Section~\ref{wkr:sec:future} is Gao's map of his
Theorem~3.5, written $G_{\mathrm{Gao}}$ here because $G(t)$ below is the
coefficient of $v$ in $r$ (after Theorem~\ref{wkr:thm:triangular} and in
Section~\ref{wkr:sec:future}), and its \emph{degree four is the geometric degree}, the
number of preimages of a generic point, the sense of ``degree'' defined
in his abstract. Its three components have ordinary degrees $4$, $11$ and
$12$ and its Jacobian determinant is $2$: this was read from his
Section~3.5 at the write and checked there with SymPy~1.14.0 from his
printed formulas (component degrees $4,11,12$ and $\det J_{G_{\mathrm{Gao}}}=2$ as
polynomial identities). So $G_{\mathrm{Gao}}$ is \emph{not} a three-dimensional Keller
map of lower ordinary degree than seven. The sentence ``our forced degree
seven is compatible with examples outside this class of lower degree''
remains true as a compatibility statement, but $G_{\mathrm{Gao}}$ is not such an example;
for Alp\"oge's map Gao's Theorem~3.3 gives geometric degree three. The
project's framing is unaffected: its notes say that the finite searches
``do not prove that ordinary degree seven or sixteen monomials is globally
minimal in dimension three'' and that a simpler three-variable map would
have degree at least four and, below seven, would break the weighted
symmetry (\rpath{Algebra/JacobianConjecture/Research/README.md}, lines
12--15 and 262--269). Whether such a map exists is not settled by any
source cited here.

\section{Weighted coordinates and the exact determinant reduction}
\label{wkr:sec:coordinates}

Throughout the classification, $K$ is an arbitrary field of characteristic
zero. All polynomial equalities are formal equalities over $K$.

\begin{definition}
A polynomial map $F:\A^3_K\to\A^3_K$ is a \emph{Keller map} if its formal
Jacobian determinant is a nonzero element of $K$. We use the source action
\[
(x,y,z)\longmapsto(\rho^{-1}x,\rho y,\rho^2z)
\]
and the target action with weights $(2,1,-1)$. Equivariance means that the
three coordinate polynomials are weighted homogeneous of weights $2,1,-1$,
respectively.
\end{definition}

\subsection{Invariant numerators}

Set
\begin{equation}\label{wkr:eq:invariants}
t=xy,\qquad v=x^2z.
\end{equation}
Both are invariant under the source action. Every equivariant polynomial
map has a unique expression in the localization at $x$ of the form
\begin{equation}\label{wkr:eq:lift}
F=\mathcal L(p,q,r)
 :=\left(\frac{p(t,v)}{x^2},\frac{q(t,v)}x,xr(t,v)\right),
\qquad p,q,r\in K[t,v].
\end{equation}
For example, a monomial $x^iy^jz^k$ of weight $2$ satisfies
$i=j+2k-2$; multiplying it by $x^2$ gives $t^jv^k$.
The corresponding formulas for weights $1$ and $-1$ are
$i=j+2k-1$ and $i=j+2k+1$. They prove existence and uniqueness of the
numerators monomial by monomial.

For general numerators, polynomiality requires that every monomial $t^iv^j$
in $p$ satisfy $i+2j\ge2$, and every such monomial in $q$ satisfy
$i+2j\ge1$. In the class studied here, $\deg_vp,\deg_vq\le1$, so these
conditions reduce exactly to
\begin{equation}\label{wkr:eq:polynomiality}
p(0,0)=p_t(0,0)=q(0,0)=0.
\end{equation}
There is no additional restriction on $r$ for polynomiality. The apparent
fractions in \eqref{wkr:eq:lift} are therefore polynomial expressions whenever
\eqref{wkr:eq:polynomiality} holds in our class.

\begin{proposition}[Weighted determinant identity]\label{wkr:prop:jacobian}
For a polynomial lift \eqref{wkr:eq:lift},
\begin{equation}\label{wkr:eq:J}
\det J_F=\J(p,q,r):=
\det\begin{pmatrix}
-2p&p_t&p_v\\
-q&q_t&q_v\\
r&r_t&r_v
\end{pmatrix},
\end{equation}
where the right side is evaluated at $(t,v)=(xy,x^2z)$.
\end{proposition}
\begin{proof}
Compute in coordinates $(x,t,v)$ on $x\ne0$. The coordinate change
$(x,y,z)\mapsto(x,xy,x^2z)$ has determinant $x^3$. The derivative matrix of
$(x^{-2}p,x^{-1}q,xr)$ with respect to $(x,t,v)$ has rows
\[
(-2x^{-3}p,x^{-2}p_t,x^{-2}p_v),\quad
(-x^{-2}q,x^{-1}q_t,x^{-1}q_v),\quad
(r,xr_t,xr_v).
\]
Factoring $x^{-2},x^{-1},x$ from its rows and $x^{-1}$ from its first column
leaves the matrix in \eqref{wkr:eq:J} and contributes $x^{-3}$. The chain rule
cancels this with $x^3$. Localization is injective on $K[x,y,z]$, so the
identity holds as a polynomial identity, including on $x=0$.
\end{proof}
\wtag{} This identity, with this proof, and the normalization below are
already in the project's research notes (Section~\ref{wkr:sub:project}).

\subsection{Normalization is exhaustive}

Suppose that \eqref{wkr:eq:polynomiality} holds and $\J$ is a nonzero constant.
Evaluating at $(t,v)=(0,0)$ gives
\begin{equation}\label{wkr:eq:originJ}
\J(0,0)=-p_v(0,0)q_t(0,0)r(0,0).
\end{equation}
All three factors are therefore nonzero. Dividing the three target
coordinates by these factors normalizes
\begin{equation}\label{wkr:eq:normalization}
p_v(0,0)=q_t(0,0)=r(0,0)=1.
\end{equation}
For a Keller map, the resulting constant determinant is $-1$.
This operation preserves the structural restrictions on $p,q,r$, ordinary
degrees, support sizes, and invertibility or noninvertibility. Thus a
classification in this normalization classifies the unnormalized maps up
to independent nonzero target scalings. We can write
\begin{equation}\label{wkr:eq:r}
r=1+\beta t+\gamma v.
\end{equation}
The word \emph{normalized} below always includes
\eqref{wkr:eq:polynomiality} and \eqref{wkr:eq:normalization}; it does not implicitly
assume $\beta=\gamma=1$.

\section{The complete all-degree classification}
\label{wkr:sec:main}

Define the following polynomials over $\Q[a,b,t,v]$:
\begin{subequations}\label{wkr:eq:family}
\begin{align}
p_{a,b}={}&v-\frac{8a}{9}t^2-2ab\,tv
 +\frac{28a^2b}{27}t^3+\frac{4a^2b^2}{3}t^2v \notag\\
 &\hspace{1.2em}-\frac{8a^3b^2}{27}t^4-\frac{8a^3b^3}{27}t^3v,
 \label{wkr:eq:familyp}\\
q_{a,b}={}&t+9bv-8ab\,t^2-12ab^2tv
 +4a^2b^2t^3+4a^2b^3t^2v,\label{wkr:eq:familyq}\\
r_{a,b}={}&1+ab\,t+ab^2v.\label{wkr:eq:familyr}
\end{align}
\end{subequations}
Let $F_{a,b}=\mathcal L(p_{a,b},q_{a,b},r_{a,b})$.

\begin{theorem}[All-degree classification]\label{wkr:thm:classification}
Let $p,q,r\in K[t,v]$ satisfy $\deg_vp,\deg_vq\le1$, let $r$ be affine in
$(t,v)$, and suppose that their lift is polynomial and normalized as in
\eqref{wkr:eq:normalization}. Then $\J(p,q,r)=-1$ if and only if exactly one
of the following descriptions applies.
\begin{enumerate}[label=\textup{(\Alph*)}]
\item\label{wkr:case:tame} $r=1$ and, for unique $\lambda\in K$ and
$B\in t^2K[t]$,
\begin{equation}\label{wkr:eq:tamepq}
p=v+B(t),\qquad q=t+\lambda\bigl(v+B(t)\bigr).
\end{equation}
Its lift is a tame polynomial automorphism.
\item\label{wkr:case:nonconstant} $r$ is nonconstant, both $\beta$ and $\gamma$
in \eqref{wkr:eq:r} are nonzero, and
\begin{equation}\label{wkr:eq:abrecover}
(p,q,r)=(p_{a,b},q_{a,b},r_{a,b}),\qquad
 a=\frac{\beta^2}{\gamma},\quad b=\frac{\gamma}{\beta}.
\end{equation}
Its lift is noninjective. Its coordinate degrees are $(7,6,4)$, and its
coordinate support sizes are $(7,6,3)$.
\end{enumerate}
In particular, there are no normalized Keller maps in this class with
exactly one of $\beta,\gamma$ nonzero. No ordinary degree bound is assumed.
\end{theorem}

We prove the classification in Sections~\ref{wkr:sec:lemmas}--\ref{wkr:sec:proof}.
The invertibility, collisions, and support assertions are proved in
Section~\ref{wkr:sec:consequences}. The forward direction is the substantive
rigidity statement; the reverse direction follows either from the formulas
there or by substitution into \eqref{wkr:eq:J}.

\subsection{The normalized representative}

Put $F_*=F_{1,1}$. Its numerators are
\begin{align*}
p_*&=v-\frac89t^2-2tv+\frac{28}{27}t^3
       +\frac43t^2v-\frac8{27}t^4-\frac8{27}t^3v,\\
q_*&=t+9v-8t^2-12tv+4t^3+4t^2v,\qquad r_*=1+t+v.
\end{align*}
It is exactly a normalization of \eqref{wkr:eq:original}:
\begin{equation}\label{wkr:eq:originalnormalization}
F_*=\diag\left(-\frac12,-\frac32,\frac12\right)
\circ F_0\circ\diag\left(1,-\frac23,-2\right).
\end{equation}
Thus the nonconstant part of the classification identifies a known
representative rather than producing an unrelated counterexample.

\section{One-variable rigidity lemmas}
\label{wkr:sec:lemmas}

The decisive step is that constant weighted Wronskians control degree even
when arbitrarily high-degree polynomials are allowed initially.

\begin{lemma}[A degree bound from a constant weighted Wronskian]
\label{wkr:lem:wronskian}
Let $B,D\in K[t]$ satisfy
\begin{equation}\label{wkr:eq:wronskian}
DB'-2BD'=\kappa,\qquad \kappa\in K^\times.
\end{equation}
Then $D\ne0$, $\deg D\le1$, and $\deg B\le2$. More precisely, either $D$
is a nonzero constant and $B$ has degree one, or $D$ has degree one and
$B=\eta D^2+\mu$ for some $\eta,\mu\in K$, with
$-2\mu D'=\kappa$.
\end{lemma}
\begin{proof}
Neither $B$ nor $D$ can be zero. If $D$ is constant, \eqref{wkr:eq:wronskian}
gives $B'=\kappa/D$, so $B$ has degree one.

Now put $n=\deg D\ge1$ and $m=\deg B$. Unless $m=2n$, the leading term
on the left has degree $m+n-1$ and nonzero coefficient proportional to
$m-2n$. This uses characteristic zero. If $m>2n$, it contradicts the
constant right-hand side. If $m=2n$, subtract from $B$ the unique scalar
multiple of $D^2$ that cancels its leading coefficient. This changes
neither side of \eqref{wkr:eq:wronskian}. The resulting polynomial $B_0$ is
nonzero and has degree $m_0<2n$. If initially $m<2n$, just take $B_0=B$.

The leading term calculation for $B_0$ now gives $m_0+n-1=0$. As
$m_0\ge0$ and $n\ge1$, we obtain $n=1$ and $m_0=0$. Hence
$B=\eta D^2+\mu$, and substitution gives $-2\mu D'=\kappa$.
\end{proof}

\begin{lemma}[A cubic--square relation]\label{wkr:lem:powers}
Suppose $A,C\in K[t]$ are nonzero and
\begin{equation}\label{wkr:eq:powers}
2CA'-3AC'=0,\qquad A(0)=1.
\end{equation}
Then there are unique $E\in K[t]$ with $E(0)=1$ and $c\in K^\times$ such
that
\begin{equation}\label{wkr:eq:EC}
A=E^3,\qquad C=cE^2.
\end{equation}
\end{lemma}
\begin{proof}
The derivative of $A^2/C^3$ is zero. A rational function over a
characteristic-zero field with derivative zero lies in $K$: writing it in
coprime numerator--denominator form, the identity $f'g-fg'=0$ forces each
nonconstant numerator or denominator to divide its own derivative, which
is impossible by degree. Thus $A^2/C^3\in K^\times$.

For every irreducible polynomial $\pi$, unique factorization gives
$2\operatorname{ord}_\pi A=3\operatorname{ord}_\pi C$. Consequently the
multiplicities in $A$ are multiples of three and those in $C$ are the
corresponding multiples of two. Because $A(0)=1$, none of the common
factors vanishes at zero. Normalize their product to have value one there.
The resulting polynomial $E$ satisfies $A=E^3$ and $C=C(0)E^2$.
No extraction of a scalar root or extension of $K$ is necessary.
Uniqueness follows from the factorization and the value at zero.
\end{proof}

\begin{lemma}[The remaining common factor is linear]\label{wkr:lem:Edegree}
Suppose that $A(0)=1$, $B(0)=-1$, and
\begin{align}
DB'-2BD'&=-1,\label{wkr:eq:E0}\\
DA'+2CB'-3AD'-2BC'&=0,\label{wkr:eq:E1}\\
2CA'-3AC'&=0.\label{wkr:eq:E2}
\end{align}
If $C\ne0$, then in \eqref{wkr:eq:EC} the polynomial $E$ has degree at most
one.
\end{lemma}
\begin{proof}
Lemma~\ref{wkr:lem:wronskian} bounds $\deg D\le1$ and $\deg B\le2$.
Lemma~\ref{wkr:lem:powers} gives $A=E^3$, $C=cE^2$. Suppose $m=\deg E\ge2$.

If $D$ has degree one, the leading coefficient of $DA'-3AD'$ in degree
$3m$ is $3(m-1)$ times the leading coefficient of $D$ times the cube of
that of $E$. It is nonzero. The other part $2CB'-2BC'$ of \eqref{wkr:eq:E1}
has degree at most $2m+1<3m$, a contradiction.

If $D$ is constant, it is nonzero and $B$ has degree one. The term $DA'$
has degree $3m-1$, whereas $2CB'-2BC'$ has degree at most $2m$.
Again the leading term cannot cancel. Thus $\deg E\le1$.
\end{proof}

\begin{remark}
The lemmas address the possible high-degree cancellations explicitly. A
coefficient search stopped at any fixed degree would not establish these
bounds. Characteristic zero is essential to the leading coefficient
arguments and to the assertion about rational functions with zero derivative.
\end{remark}

\section{Proof of the classification, including exceptional cases}
\label{wkr:sec:proof}

\subsection{\texorpdfstring{The case $\gamma\ne0$: a moving coordinate}{The case gamma nonzero: a moving coordinate}}

First reduce $r=1+\beta t+\gamma v$ to coefficient one on $v$. Replace
\begin{equation}\label{wkr:eq:gammachange}
\widetilde p(t,v)=\gamma p(t,v/\gamma),\quad
\widetilde q(t,v)=q(t,v/\gamma),\quad
\widetilde r(t,v)=r(t,v/\gamma).
\end{equation}
The determinant and normalization are unchanged, while
$\widetilde r=1+\beta t+v$. We temporarily omit the tildes.

Set $s=1+\beta t+v$ and write uniquely
\begin{equation}\label{wkr:eq:ABCD}
p=A(t)s+B(t),\qquad q=C(t)s+D(t),\qquad r=s.
\end{equation}
Substitution into \eqref{wkr:eq:J} gives
\begin{equation}\label{wkr:eq:threeidentities}
\J=E_0+E_1s+E_2s^2,
\end{equation}
where the expressions $E_0,E_1,E_2$ are the left sides of
\eqref{wkr:eq:E0}--\eqref{wkr:eq:E2}. For clarity, differentiation at fixed $v$
gives $p_t=A's+A\beta+B'$, $q_t=C's+C\beta+D'$, and $r_t=\beta$.
Expanding the determinant cancels the $\beta$ terms and yields
\eqref{wkr:eq:threeidentities}. Since $(t,s)$ are independent polynomial
coordinates, the Keller equation is exactly
\eqref{wkr:eq:E0}--\eqref{wkr:eq:E2}.

The boundary conditions become
\begin{equation}\label{wkr:eq:boundaryABCD}
\begin{gathered}
A(0)=1,\quad B(0)=-1,\quad A'(0)+\beta+B'(0)=0,\\
C(0)+D(0)=0,\qquad C'(0)+\beta C(0)+D'(0)=1.
\end{gathered}
\end{equation}
We must not divide by $C$ before excluding $C=0$. If $C$ were zero, the
fourth condition would give $D(0)=0$. Equation~\eqref{wkr:eq:E0} at zero would
then give $D'(0)=-1/2$, while \eqref{wkr:eq:E1} at zero would give
$-3D'(0)=0$. This is impossible.

The rigidity lemmas now apply. Write
\begin{equation}\label{wkr:eq:linearforms}
\begin{gathered}
E=1+kt,\quad A=E^3,\quad C=cE^2,\qquad c\ne0,\\
D=dt-c,\qquad B=-1+b_1t+b_2t^2.
\end{gathered}
\end{equation}
The first derivative condition in \eqref{wkr:eq:boundaryABCD} gives
\begin{equation}\label{wkr:eq:b1}
b_1=-\beta-3k.
\end{equation}
After dividing \eqref{wkr:eq:E1} by the nonzero polynomial $E$, it reads
\begin{equation}\label{wkr:eq:keylinear}
2c(EB'-2kB)=3(kc+d)E.
\end{equation}
Its constant and linear coefficients, together with \eqref{wkr:eq:b1}, give
\begin{equation}\label{wkr:eq:bd}
b_2=-\beta k-2k^2,\qquad
\frac dc=-\frac{2\beta+5k}{3}.
\end{equation}
The final condition in \eqref{wkr:eq:boundaryABCD} is
$c(\beta+2k)+d=1$, which therefore simplifies to
\begin{equation}\label{wkr:eq:ck}
c(\beta+k)=3.
\end{equation}
The coefficient of $t$ in \eqref{wkr:eq:E0} is $-db_1-2cb_2$. Its vanishing,
using \eqref{wkr:eq:b1} and \eqref{wkr:eq:bd}, is
\[
\frac c3(\beta+k)(2\beta+3k)=0.
\]
Equation~\eqref{wkr:eq:ck} permits cancellation of the first factors. Hence
\begin{equation}\label{wkr:eq:endpoint}
k=-\frac{2\beta}{3},\qquad
\beta\ne0,\qquad c=\frac9\beta,\qquad d=4,
\qquad b_1=\beta,\qquad b_2=-\frac{2\beta^2}{9}.
\end{equation}
The remaining constant coefficient of \eqref{wkr:eq:E0} is indeed $-1$:
$-cb_1+2d=-9+8=-1$. All three equations hold for these values. We have proved
uniqueness and existence, including the exclusion of $\beta=0$.

More compactly, before undoing \eqref{wkr:eq:gammachange}, the answer is
\begin{equation}\label{wkr:eq:compactendpoint}
\begin{gathered}
E=1-\frac{2\beta}{3}t,\qquad s=1+\beta t+v,\\
p=E^3s-1+\beta t-\frac{2\beta^2}{9}t^2,\qquad
q=\frac9\beta E^2s+4t-\frac9\beta.
\end{gathered}
\end{equation}
Replace $v$ by $\gamma v$ and divide $p$ by $\gamma$. Expansion yields
\eqref{wkr:eq:family} with $a=\beta^2/\gamma$ and $b=\gamma/\beta$.
This proves the nonconstant branch whenever $\gamma\ne0$.

\subsection{\texorpdfstring{The case $\gamma=0$, $\beta\ne0$: no solutions}{The case gamma zero, beta nonzero: no solutions}}

Write $p=A(t)v+B(t)$, $q=C(t)v+D(t)$, and $r=1+\beta t$. The normalization
implies $A(0)=1$ and $B(0)=B'(0)=D(0)=0$. The coefficient of $v$ in
\eqref{wkr:eq:J} is
\begin{equation}\label{wkr:eq:gamma0v}
\beta AC+r(A'C-AC')=0.
\end{equation}
If $C\ne0$, this says that $(rA/C)'=0$, so $C=\lambda rA$ for some
$\lambda\in K^\times$. When $C=0$, the same expression holds with
$\lambda=0$. The constant coefficient of the determinant then becomes
\begin{equation}\label{wkr:eq:gamma0const}
A\left(\lambda(r^2B)'-(rD)'\right)=-1.
\end{equation}
A product of two polynomials is a nonzero constant only when both are
constants. Thus $A=1$. Integrating \eqref{wkr:eq:gamma0const} and using
$B(0)=D(0)=0$ gives
\begin{equation}\label{wkr:eq:impossible}
\lambda r^2B-rD=-t.
\end{equation}
Evaluate at $t=-1/\beta$. The left side is zero and the right side is
$1/\beta\ne0$. This contradiction excludes the entire case, at every degree.

\subsection{\texorpdfstring{The case $r=1$: the full tame branch}{The case r=1: the full tame branch}}

Again write $p=Av+B$ and $q=Cv+D$. The coefficient of $v$ in the determinant
is $A'C-AC'=0$. Since $A(0)=1$, it follows that $C=\lambda A$ for some
$\lambda\in K$, allowing $C=0$. The constant coefficient is
\[
A(\lambda B'-D')=-1.
\]
Consequently $A=1$ and $D=t+\lambda B$, using $B(0)=D(0)=0$.
Polynomiality and normalization give $B(0)=B'(0)=0$, hence $B\in t^2K[t]$.
This is precisely \eqref{wkr:eq:tamepq}. Conversely, these formulas give
$\J=-1$ directly. Together with the preceding two cases, this completes
the algebraic classification in Theorem~\ref{wkr:thm:classification}.

\section{Invertibility, diagonal equivalence, and sharp support}
\label{wkr:sec:consequences}

\subsection{An inverse for every constant-third-numerator map}

Write $B(t)=t^2H(t)$. The lift of \eqref{wkr:eq:tamepq} is
\begin{equation}\label{wkr:eq:tameF}
F(x,y,z)=\bigl(P,\ y+\lambda xP,\ x\bigr),\qquad
P=z+y^2H(xy).
\end{equation}
For output coordinates $(U,V,W)$, define
\begin{equation}\label{wkr:eq:tameinverse}
 x=W,\qquad y=V-\lambda WU,\qquad
 z=U-(V-\lambda WU)^2H\bigl(W(V-\lambda WU)\bigr).
\end{equation}
Substitution in both directions is immediate from $P=U$. This proves
polynomial invertibility. More precisely, \eqref{wkr:eq:tameF} is the
composition of the elementary shear
$z\mapsto z+y^2H(xy)$, the elementary shear $y\mapsto y+\lambda xz$ in the
updated coordinates, and the coordinate permutation
$(x,y,z)\mapsto(z,y,x)$. It is therefore tame. There is no bound on
$\deg H$: arbitrarily large degrees occur on the invertible branch.

\subsection{Every nonconstant map is in the same diagonal orbit}

For $a,b\ne0$, inspection of \eqref{wkr:eq:family} gives
\begin{equation}\label{wkr:eq:diagonal}
F_{a,b}=\diag\bigl((ab^2)^{-1},(ab)^{-1},1\bigr)
\circ F_*\circ\diag(1,ab,ab^2).
\end{equation}
The product of the two diagonal determinants is one, consistently with
$\det J_{F_{a,b}}=-1$. Thus there are two parameters before taking this
source--target equivalence, and no parameter after taking it. The scalar
parameters are uniquely recoverable from either $r$ or two coefficients:
\begin{equation}\label{wkr:eq:coeffrecover}
a=-\frac98\coeff{t^2}{p},\qquad
b=\frac19\coeff{v}{q}.
\end{equation}

\begin{corollary}[Sharp all-degree support rigidity]\label{wkr:cor:support}
Within the class of Theorem~\ref{wkr:thm:classification}, every noninvertible
map has ordinary maximum degree seven and exactly sixteen expanded
monomials. After ordering coordinates by their prescribed weights, its
degree profile is $(7,6,4)$ and its support profile is $(7,6,3)$.
\end{corollary}
\begin{proof}
Only branch~\ref{wkr:case:nonconstant} can be noninvertible. For $a,b\ne0$,
all coefficients displayed in \eqref{wkr:eq:family} are nonzero. Different
monomials in a numerator lift to different monomials. The degree of a lifted
$t^iv^j$ is $2i+3j-2$ in the first coordinate, $2i+3j-1$ in the second,
and $2i+3j+1$ in the third. The asserted profiles follow. Noninjectivity of
each map is established just below.
\end{proof}

This is stronger than saying that sixteen terms are minimal under a degree
cap: every noninvertible map in the structural class has that precise
support profile, regardless of the degree initially allowed. It is not a
statement about minimum degree or support outside the class.

\subsection{Collisions valid over the ground field}

Let $\gamma=ab^2$ and $\ell=9b$. Direct evaluation on the plane $x=0$
gives
\begin{equation}\label{wkr:eq:xzero}
F_{a,b}(0,y,z)=\left(z-\frac{8a}{9}y^2,y,0\right).
\end{equation}
At $t=0$, the numerators simplify to $p=v$, $q=\ell v$, and
$r=1+\gamma v$. Therefore
\begin{equation}\label{wkr:eq:newcollision}
F_{a,b}\left(1,0,-\frac1\gamma\right)
=F_{a,b}\left(0,-\frac\ell\gamma,\frac{71}\gamma\right)
=\left(-\frac1\gamma,-\frac\ell\gamma,0\right).
\end{equation}
To check the constant $71$, note that
$-(8a/9)\ell^2=-72\gamma$; the first coordinate of the second point's
image is therefore $(71-72)/\gamma$.

The two source points have different $x$ coordinates. This proves
noninjectivity over \emph{every} characteristic-zero ground field $K$,
without passing to an algebraic closure. It also gives a rational
certificate when $a,b\in\Q^\times$.

Equivariance propagates this to every $\rho\in K^\times$:
\begin{equation}\label{wkr:eq:collisionfamily}
\begin{aligned}
F_{a,b}\left(\rho,0,-\frac1{\gamma\rho^2}\right)
&=F_{a,b}\left(0,-\frac\ell{\gamma\rho},
                         \frac{71}{\gamma\rho^2}\right)\\
&=\left(-\frac1{\gamma\rho^2},-\frac\ell{\gamma\rho},0\right).
\end{aligned}
\end{equation}
These are explicit certificates in the coordinates of this report. They
are consequences of the classified family, not a claim that the collision
mechanism of the original map was unknown.
\wtag{} The mechanism is the project's
(\rpath{Algebra/JacobianConjecture/Research/README.md}, lines 210--224):
\eqref{wkr:eq:newcollision} is its case $t_0=0$, $v_0=-1/\gamma$, and the
second point $(0,q,p-\alpha q^2)$ with $\alpha=[t^2]p_{a,b}=-8a/9$ (the
notes' $a$) is the one displayed. For $a=b=1$ this was checked at the write by exact evaluation.

\begin{example}[The normalized rational witness]
For $a=b=1$, the points $(1,0,-1)$ and $(0,-9,71)$ have common image
$(-1,-9,0)$. The Jacobian is nonsingular at both points because it has
constant determinant $-1$. Thus local nonsingularity and global injectivity
are genuinely different conditions in the displayed family.
\end{example}

\subsection{A larger class: arbitrary dependence on the first invariant}
\label{wkr:sec:triangular}

The affine restriction on $r$ can be relaxed substantially without changing
its essential normal form. The relevant source automorphisms are
\begin{equation}\label{wkr:eq:source-shear}
T_h(x,y,z)=\bigl(x,y,z+y^2h(xy)\bigr),\qquad h\in K[t].
\end{equation}
They preserve the weights, have Jacobian one, and induce
$(t,v)\mapsto(t,v+t^2h(t))$. Their inverses are $T_{-h}$.

\begin{theorem}[All-degree triangular extension]\label{wkr:thm:triangular}
Let $p,q$ be affine in $v$ and let
\[
r=R(t)+\gamma v,\qquad R\in K[t],\quad \gamma\in K,
\]
with polynomial lift and normalized origin coefficients. In particular,
$R(0)=1$. Then the normalized lift is Keller if and only if one of the
following holds.
\begin{enumerate}[label=\textup{(\roman*)}]
\item $\gamma=0$, $R=1$, and the map is a tame map
\eqref{wkr:eq:tameF}.
\item $\gamma\ne0$, $\beta:=R'(0)\ne0$, and the map is
\begin{equation}\label{wkr:eq:triangular-normal}
F=F_{a,b}\circ T_h,\qquad
 a=\frac{\beta^2}{\gamma},\quad b=\frac{\gamma}{\beta},\quad
 h(t)=\frac{R(t)-1-\beta t}{\gamma t^2}.
\end{equation}
Here $h$ is a polynomial, uniquely determined by $r$. The map is
noninjective over $K$.
\end{enumerate}
Thus every noninvertible map in this larger class is the classified
representative up to diagonal scalings and one explicit equivariant source
shear. There is no degree bound on $R$, $p$, or $q$.
\end{theorem}
\begin{proof}
When $\gamma\ne0$, the numerator defining $h$ vanishes to order at least
two at zero. Composing the source with $T_{-h}$ replaces $v$ by
$v-t^2h(t)$ and changes the third numerator to $1+\beta t+\gamma v$.
The first two numerators remain affine in $v$; polynomiality and all
normalized origin coefficients are preserved. The source shear has
Jacobian one. Theorem~\ref{wkr:thm:classification} therefore applies: it
forces $\beta\ne0$ and identifies $F\circ T_{-h}$ with $F_{a,b}$.
This is \eqref{wkr:eq:triangular-normal}. The converse follows from the same
composition, and collision points are transported by $T_{-h}$.

When $\gamma=0$, use the proof of Section~\ref{wkr:sec:proof} with
$r=R(t)$ instead of $1+\beta t$. The coefficient of $v$ is
$r'AC+r(A'C-AC')=0$, whence $C=\lambda rA$, including $C=0$.
The constant equation is still
$A\bigl(\lambda(r^2B)'-(rD)'\bigr)=-1$.
It forces $A=1$ and, on integration with the origin conditions,
$\lambda r^2B-rD=-t$. Thus $r$ divides $t$. Since $r(0)=1$, this
forces $r$ to be constant and hence $r=1$. The already proved tame
classification finishes this case.
\end{proof}

\begin{corollary}[An exact degree spectrum]\label{wkr:cor:degreespectrum}
In the noninvertible branch of Theorem~\ref{wkr:thm:triangular}, $h=0$ gives
coordinate degrees $(7,6,4)$. If $h\ne0$ has degree $m$, the coordinate
degrees are exactly
\begin{equation}\label{wkr:eq:extendeddegrees}
(2m+8,\ 2m+7,\ 2m+5).
\end{equation}
Consequently the possible maximum ordinary degrees in this noninvertible
class are precisely $7$ and every even integer at least $8$.
\end{corollary}
\begin{proof}
The coefficients of $z$ in the three coordinates of $F_{a,b}$ have
ordinary degrees $6,5,3$, respectively. Their leading terms are nonzero
for $a,b\ne0$. The substitution $z\mapsto z+y^2h(xy)$ adds terms of
degrees $6+(2+2m)$, $5+(2+2m)$, and $3+(2+2m)$. These are all strictly
larger than the original coordinate degrees, so they cannot cancel against
an original term. Their leading coefficients are nonzero, proving
\eqref{wkr:eq:extendeddegrees}. Every $m\ge0$ occurs by taking $h=t^m$.
\end{proof}

The first invariant coefficient is the important distinction here. The
coefficient of $v$ in $r$ is still required to be a scalar, not a
nonconstant polynomial in $t$. No assertion is made for
$r=R(t)+G(t)v$ with nonconstant $G$.
[Added 29 September 2026, batch~44: when $p$ and $q$ are affine in $v$,
a nonconstant $G$ is impossible (Part~II,
Theorem~\ref{wkr:as:thm:slope}), so Theorem~\ref{wkr:thm:triangular}
classifies every such map with $r$ affine in $v$
(Theorem~\ref{wkr:as:thm:classification}).]

\section{Reduced parameter geometry and invertible degenerations}
\label{wkr:sec:geometry}

The formulas \eqref{wkr:eq:family} contain no negative powers of $a$ or $b$.
They therefore define a polynomial family on the entire parameter plane,
not just on its torus. The diagonal formula \eqref{wkr:eq:diagonal} should not
be used on the axes, where its individual scaling matrices are undefined.
The polynomial formulas are the correct extension there.

\begin{theorem}[The parameter plane and its boundary]\label{wkr:thm:plane}
In the finite coefficient space of normalized weighted maps of ordinary
degree at most seven with $p,q$ affine in $v$ and $r$ affine in $(t,v)$,
the following statements hold over $K$.
\begin{enumerate}[label=\textup{(\roman*)}]
\item The morphism $(a,b)\mapsto F_{a,b}$ is a closed embedding of
$\A^2_K$. Every member has Jacobian determinant $-1$.
\item Its torus $ab\ne0$ is exactly the normalized nonconstant-$r$ branch.
Its reduced Zariski closure is the entire embedded affine plane.
\item The invertible members of this plane are exactly its two coordinate
axes $V(ab)$. Explicitly,
\begin{align}
F_{0,b}(x,y,z)&=(z,\ y+9bxz,\ x),\label{wkr:eq:axisA}\\
F_{a,0}(x,y,z)&=\left(z-\frac{8a}{9}y^2,\ y,\ x\right).
\label{wkr:eq:axisB}
\end{align}
\item Inside this plane, the scheme-theoretic condition $r=1$ cuts out the
reduced nodal union $\Spec K[a,b]/(ab)$.
\end{enumerate}
\end{theorem}
\begin{proof}
The coefficients \eqref{wkr:eq:coeffrecover} recover $a,b$ by linear functions
on the ambient coefficient space. Every other coefficient is a polynomial
in those two functions by \eqref{wkr:eq:family}. Thus the image is the graph of
polynomial functions, a closed embedding.

The determinant is $-1$ on the torus by the classification and its explicit
solution. It is a polynomial in $a,b,t,v$, so the identity holds on the
whole plane. Equivalently, it can be checked directly from
\eqref{wkr:eq:family}. The torus is dense, and the classification identifies all
its field-valued nonconstant-$r$ points. This proves the closure statement.

On the axes, the formulas reduce to \eqref{wkr:eq:axisA} and \eqref{wkr:eq:axisB},
which are tame maps of the already classified constant branch. Off the
axes, \eqref{wkr:eq:newcollision} excludes invertibility. Finally, imposing
$r=1$ on the parameter ring imposes $(ab,ab^2)=(ab)$. The ideal $(ab)$ is
radical: it is the intersection $(a)\cap(b)$ in $K[a,b]$.
\end{proof}

\subsection{What degenerates, and what does not}

The Jacobian determinant stays $-1$ throughout this degeneration. The
ordinary degrees and the collision coordinates do not stay bounded in the
same manner. Fixing $b=1$ and letting $a$ approach zero over $\mathbb R$ or
$\mathbb C$, the coefficients tend to the automorphism $F_{0,1}$, whereas
the two source points in \eqref{wkr:eq:newcollision} have coordinates with
poles in $a$. In this explicit sense the displayed collisions escape to
infinity. No bounded limiting pair of distinct colliding points is asserted.

There are many other constant-$r$ maps, because $B\in t^2K[t]$ in
\eqref{wkr:eq:tamepq} is arbitrary. Only the two axes described above lie in the
closure of the nonconstant branch in this coefficient space. In particular,
passing from the nonconstant branch to its closure does not recover the
entire infinite-dimensional tame family.

\subsection{Rigidity over every reduced coefficient algebra}

Field-valued uniqueness can be strengthened without any degree cap on the
numerators.

\begin{corollary}[All-degree reduced rigidity]\label{wkr:cor:reduced}
Let $R$ be a reduced $\Q$-algebra. Suppose $p,q\in R[t,v]$ are affine in
$v$, satisfy the normalization and polynomiality conditions, and satisfy
$\J(p,q,1+t+v)=-1$. Then $p=p_*$ and $q=q_*$.
\end{corollary}
\begin{proof}
For every prime ideal $\mathfrak p$ of $R$, pass to the fraction field of
$R/\mathfrak p$. It is a characteristic-zero field, so
Theorem~\ref{wkr:thm:classification} forces $a=b=1$. Every coefficient of
$p-p_*$ and $q-q_*$ therefore belongs to every prime ideal of $R$, hence to
its nilradical. That nilradical is zero because $R$ is reduced.
\end{proof}

Thus fixing $r=1+t+v$ eliminates nonconstant families over reduced bases,
even when large ordinary degrees are allowed. It does \emph{not} eliminate
families over nonreduced bases. We now determine one such coefficient
scheme exactly.

\section{The normalized coefficient scheme is a double point}
\label{wkr:sec:scheme}

The preceding proof used degree arguments over fields. Nilpotent leading
coefficients need not satisfy those arguments. Scheme structure must
therefore be analyzed rather than inferred from the unique field-valued
solution.

\subsection{The precise finite coefficient problem}

Fix $r=1+t+v$ and ordinary degree at most seven. The most general normalized
numerators in our class are
\begin{align}
p={}&v+p_{20}t^2+p_{11}tv+p_{30}t^3+p_{21}t^2v
       +p_{40}t^4+p_{31}t^3v,\label{wkr:eq:scheme-p}\\
q={}&t+q_{01}v+q_{20}t^2+q_{11}tv+q_{30}t^3
       +q_{21}t^2v+q_{40}t^4.\label{wkr:eq:scheme-q}
\end{align}
The degree formulas in the proof of Corollary~\ref{wkr:cor:support} show that
this list is exhaustive. In particular, the $q_{40}$ term is allowed and
must not be set to zero beforehand.

Let the coefficient vector, in the following fixed order, be
\begin{equation}\label{wkr:eq:coordorder}
c=(p_{20},p_{11},p_{30},p_{21},p_{40},p_{31},
       q_{01},q_{20},q_{11},q_{30},q_{21},q_{40}).
\end{equation}
Let $I\subset\Q[c]$ be generated by all coefficients in $t,v$ of
$\J(p,q,1+t+v)+1$. There are eighteen nonzero coefficient polynomials.
Their complete expressions, ordering, and derivative matrix are supplied
in \texttt{verification.json}; the article fixes enough data below to
reconstruct and check the certificates independently.

\begin{theorem}[Exact double-point presentation]\label{wkr:thm:double}
The full coordinate algebra of this coefficient problem is
\begin{equation}\label{wkr:eq:dualnumbers}
\Q[c]/I\ \simeq\ \Q[\eps]/(\eps^2).
\end{equation}
Under this isomorphism $\eps=q_{21}-4$ and
\begin{equation}\label{wkr:eq:pointplus}
c=\cstar+u\eps,
\end{equation}
where
\begin{equation}\label{wkr:eq:cstar}
\cstar=\left(-\frac89,-2,\frac{28}{27},\frac43,-\frac8{27},-\frac8{27},
                     9,-8,-12,4,4,0\right)
\end{equation}
and
\begin{equation}\label{wkr:eq:u}
u=\frac1{144}(-7,-18,18,24,-8,-8,162,-198,-324,144,144,0).
\end{equation}
In particular, the entire ideal, not merely its radical, is generated by
$\eps^2$ and the eleven linear relations obtained from
\eqref{wkr:eq:pointplus} by omitting the tautological $q_{21}$ relation.
\end{theorem}

The proof has three components: uniqueness of geometric points, a rank
certificate for the tangent space, and a certificate excluding a
second-order lift. None requires a Gr\"obner-basis computation.

\subsection{Tangent space and its exact rank certificate}

Write the coefficient equations as $f_1,\ldots,f_{18}$ in the order of the
monomials
\begin{equation}\label{wkr:eq:monomialorder}
\begin{gathered}
t^7,\ t^6v,\ t^6,\ t^5v,\ t^5,\ t^4v,\ t^4,\ t^3v^2,\ t^3v,\ t^3,\\
t^2v^2,\ t^2v,\ t^2,\ tv^2,\ tv,\ t,\ v^2,\ v.
\end{gathered}
\end{equation}
Let
\begin{equation}\label{wkr:eq:M}
M=\left(\frac{\partial f_i}{\partial c_j}(\cstar)\right)_{i,j}.
\end{equation}
Direct rational calculation gives $Mu=0$. Delete the column corresponding
to $q_{21}$ and select the rows corresponding to
\begin{equation}\label{wkr:eq:minorrowmonomials}
t^7,\ t^6,\ t^5,\ t^4v,\ t^4,\ t^3,\ t^2v,\ t^2,\ tv^2,\ t,\ v.
\end{equation}
The determinant of that $11\times11$ submatrix is
\begin{equation}\label{wkr:eq:minorvalue}
-18\,874\,368\ne0.
\end{equation}
Thus $\operatorname{rank}M=11$ and $\ker M=\Q u$. This proves that the
scheme has a one-dimensional Zariski tangent space at its sole geometric
point. It does not yet determine the length of the point.

Because $f_i(\cstar)=0$ and $Mu=0$, the assignment
$c=\cstar+u\eps$ solves the equations over $\Q[\eps]/(\eps^2)$.
Since $u_{q_{21}}=1$, the resulting map from the coefficient algebra onto
the dual numbers is surjective. In particular, the tangent direction is
not merely a redundant parametrization of a reduced point.

\subsection{A nonzero obstruction at second order}

Define the quadratic residual vector
\begin{equation}\label{wkr:eq:H}
H=\bigl([\eps^2]f_i(\cstar+u\eps)\bigr)_{i=1}^{18}.
\end{equation}
Consider a possible lift to $\Q[\eps]/(\eps^3)$ of the form
$c=\cstar+u\eps+w\eps^2$. Its second-order equation is
\begin{equation}\label{wkr:eq:lift-obstruction}
Mw+H=0.
\end{equation}
There is an explicit row functional
\begin{equation}\label{wkr:eq:ell}
L=(-205,0,-46,0,-40,24,-8,0,4,0,0,0,0,0,0,0,0,0)
\end{equation}
satisfying
\begin{equation}\label{wkr:eq:obstructionpairing}
LM=0,\qquad LH=-\frac1{72}.
\end{equation}
Applying $L$ to \eqref{wkr:eq:lift-obstruction} produces a contradiction.
For a nonzero scalar multiple $\eta u$, the obstruction pairing is
$-\eta^2/72$, which is again nonzero over every characteristic-zero field.
There is therefore no second-order lift of any nonzero first-order tangent
direction. The matrix, vector, and pairing in
\eqref{wkr:eq:M}--\eqref{wkr:eq:obstructionpairing} are checked exactly by the
included script.

\subsection{From tangent obstruction to the entire scheme}

We record the small algebraic fact that turns these certificates into a
scheme-theoretic theorem.

\begin{lemma}[An isolated point with one obstructed tangent]
\label{wkr:lem:doublecriterion}
Let $k$ be algebraically closed and let $A$ be a finitely generated
$k$-algebra with exactly one geometric point and residue field $k$.
Suppose its tangent space has dimension one and no nonzero tangent
homomorphism $A\to k[\eps]/(\eps^2)$ extends to
$A\to k[\eps]/(\eps^3)$. Then $A\simeq k[z]/(z^2)$.
\end{lemma}
\begin{proof}
The Nullstellensatz implies that the radical of zero is the unique
maximal ideal $\mathfrak m$ \cite{nullstellensatz}. The algebra is
Noetherian. A finitely generated ideal consisting of nilpotent elements
is nilpotent: if its generators have nilpotence exponents $n_i$, every
product of more than $\sum_i(n_i-1)$ generators is zero. Thus some
power of $\mathfrak m$ is zero, and $A$ is a finite-dimensional local
$k$-algebra.

The cotangent space $\mathfrak m/\mathfrak m^2$ has dimension one.
Choose $z\in\mathfrak m$ whose residue spans it. Nakayama's lemma gives
$\mathfrak m=(z)$ \cite{nakayama}. Repeatedly writing an element as a
scalar plus $z$ times another element, and using nilpotence, shows that
$A$ is generated by $z$ over $k$. Since $k[Z]$ is a principal ideal
domain and the only root of the defining ideal is zero, it follows that
$A\simeq k[Z]/(Z^N)$ for some $N$. The nonzero cotangent space forces
$N\ge2$. If $N\ge3$, sending $Z$ to $\eps$ defines a homomorphism to
$k[\eps]/(\eps^3)$ with nonzero first-order part, contrary to the
hypothesis. Hence $N=2$.
\end{proof}

\begin{proof}[Proof of Theorem~\ref{wkr:thm:double}]
Work first over $\overline\Q$. The classification with $\beta=\gamma=1$
forces $a=b=1$, so the coefficient scheme has exactly one geometric point,
$\cstar$. Equations~\eqref{wkr:eq:minorvalue} and $Mu=0$ give tangent
dimension one, and \eqref{wkr:eq:obstructionpairing} excludes every nonzero
second-order lift. Lemma~\ref{wkr:lem:doublecriterion} gives a two-dimensional
coordinate algebra over $\overline\Q$.

Already over $\Q$, substitution $c=\cstar+u\eps$ defines a surjection
\[
\Q[c]/I\ \longrightarrow\ \Q[\eps]/(\eps^2).
\]
After extension to $\overline\Q$, this is a surjection between
vector spaces of dimension two and is an isomorphism. Field extension is
faithful, so its kernel over $\Q$ was zero. This proves
\eqref{wkr:eq:dualnumbers} with the specified coordinates.

The kernel of the explicit substitution map from the polynomial ring
$\Q[c]$ is the ideal generated by $\eps^2$ and the eleven affine-linear
relations $c_j-(\cstar)_j-u_j\eps$. The isomorphism identifies this kernel
with $I$, proving the exact ideal assertion.
\end{proof}

\subsection{The universal meaning of the answer}

\begin{corollary}[Solutions over arbitrary coefficient algebras]
\label{wkr:cor:functor}
For every $\Q$-algebra $R$, the solutions of
\eqref{wkr:eq:scheme-p}--\eqref{wkr:eq:scheme-q} with $\J=-1$ are in bijection with
\[
\{e\in R:e^2=0\}.
\]
The corresponding solution is $c=\cstar+ue$, and $e=q_{21}-4$.
\end{corollary}
\begin{proof}
Apply $\operatorname{Hom}_{\Q\text{-alg}}(-,R)$ to
\eqref{wkr:eq:dualnumbers}.
\end{proof}

In particular, every solution over a reduced ring is the constant one, but
a coefficient ring containing a nonzero square-zero element supports a
nontrivial infinitesimal deformation. The parameter is obstructed rather
than the initial segment of a moving curve of normalized field-valued
solutions. The square seen in the repository's finite coefficient
certificate is therefore genuine scheme structure, not just an artifact
of one Gr\"obner-basis presentation.

\subsection{The variable-parameter open scheme}

The nonconstant coefficient parameters can also be restored without
losing control of nilpotents. Let $\mathcal X^{\times}$ denote the
coefficient scheme obtained by using \eqref{wkr:eq:scheme-p}--\eqref{wkr:eq:scheme-q},
putting $r=1+\beta t+\gamma v$, and inverting $\beta\gamma$.

\begin{theorem}[Torus times the double point]\label{wkr:thm:torusdouble}
There is an explicit isomorphism
\begin{equation}\label{wkr:eq:torusdouble}
\mathcal X^{\times}\simeq
\Spec\Q[\beta^{\pm1},\gamma^{\pm1},\eps]/(\eps^2).
\end{equation}
It is an isomorphism of full schemes, not just of reduced varieties.
The square-zero coordinate can be taken to be
$\eps=q_{21}/(\beta\gamma)-4$.
\end{theorem}
\begin{proof}
Over any $\Q$-algebra where $\beta,\gamma$ are units, define
\begin{equation}\label{wkr:eq:scheme-normalization}
\bar p(t,v)=\gamma p(t/\beta,v/\gamma),\qquad
\bar q(t,v)=\beta q(t/\beta,v/\gamma).
\end{equation}
The third numerator becomes $1+t+v$. The two normalized origin
coefficients remain one. These substitutions preserve the allowed
monomial supports and ordinary degree bound, and have polynomial inverse
formulas over the Laurent coefficient ring. They preserve the determinant
condition: the row scalings contribute $\beta\gamma$ and the derivative
column scalings contribute $(\beta\gamma)^{-1}$.

Thus \eqref{wkr:eq:scheme-normalization} identifies the coefficient functor
with a choice of two units and a solution of the fixed normalized scheme.
Apply Theorem~\ref{wkr:thm:double} to that fixed scheme. The coefficient of
$t^2v$ in $\bar q$ is $q_{21}/(\beta\gamma)$, giving the stated
square-zero coordinate. All formulas are valid with nilpotent
coefficients, so this proves the scheme isomorphism.
\end{proof}

At every geometric point of this open scheme the reduced space has
dimension two and the tangent space has dimension three: two tangent
directions come from the torus and one from its square-zero thickening.
The remaining boundary problem is how this structure meets the axes,
where the Laurent normalization is no longer available.

\begin{remark}[The degree boundary matters]
Corollary~\ref{wkr:cor:reduced} is an all-degree statement. In contrast,
Theorem~\ref{wkr:thm:double} fixes the ordinary degree-seven coefficient space.
We have not proved that allowing higher-degree coefficients over
nonreduced rings introduces no extra tangent directions or higher-order
nilpotents. The field-valued classification alone does not imply that
stronger assertion.
\end{remark}

\section{Decision procedure, reproducibility, and formalization}
\label{wkr:sec:verification}

\subsection{A classification-based recognition procedure}

The theorem gives more than a determinant identity. For rational input
polynomials in the structural class, it yields a direct recognition
procedure for the Keller condition, together with an inverse or collision
certificate.

First check the prescribed weights monomial by monomial and extract the
unique invariant numerators as in Section~\ref{wkr:sec:coordinates}. Check
$\deg_vp,\deg_vq\le1$ and that $r$ is affine. Inputs outside these
restrictions are \emph{outside the procedure's scope}; they are not
rejected as non-Keller maps.

Within this scope, compute $p_v(0),q_t(0),r(0)$. If one is zero, the map
is not Keller by \eqref{wkr:eq:originJ}. Otherwise normalize the target
coordinates. Write $r=1+\beta t+\gamma v$. If $\beta=\gamma=0$, take
$B=p-v$ and $\lambda=\coeff{v}{q}$, and test the exact identities
\eqref{wkr:eq:tamepq} with $B\in t^2K[t]$. If exactly one of $\beta,\gamma$
is nonzero, the map is not Keller. If both are nonzero, recover $a,b$ from
\eqref{wkr:eq:abrecover} and compare with the fixed formula
\eqref{wkr:eq:family}. The classification proves both soundness and completeness
of these comparisons inside the stated scope.

For a positive result, return either the inverse
\eqref{wkr:eq:tameinverse} or the pair of source points in
\eqref{wkr:eq:newcollision}. To undo target normalization, compose an inverse
with the inverse target scaling; collision source points remain collision
source points, with their common target rescaled. This avoids a search over
all possible inverses or collision pairs.

The larger class of Theorem~\ref{wkr:thm:triangular} is recognized by first
computing its unique source shear. For $\gamma\ne0$, compute $h$ from
\eqref{wkr:eq:triangular-normal}, compose the source with $T_{-h}$, and apply
the affine-$r$ procedure. For $\gamma=0$, a nonconstant $R(t)$ is excluded
by that theorem. Inputs with nonconstant $v$-coefficient in $r$ remain
outside the classified scope.
[Added 29 September 2026, batch~44: they are now classified, and
rejected, by Part~II's Theorem~\ref{wkr:as:thm:slope}; the extended
procedure is in Section~\ref{wkr:as:sec:algorithm}.]

For canonical sparse input, reading and comparing the coefficient lists
requires a number of field operations linear in the input support size,
apart from the cost of canonicalizing the input. This is an arithmetic
operation statement, not a claim of linear bit complexity for arbitrarily
large rational coefficients. An implementation should retain that distinction.

\subsection{What the companion script checks}

The file \texttt{verify\_results.py} (\wtag{} shipped as
\rpath{code/verify_results.py}, with its record in \rpath{data/}) works in
exact rational and symbolic arithmetic. It reconstructs the equations rather than trusting a stored
matrix. Its thirteen groups of checks cover:
\begin{enumerate}[label=\arabic*.]
\item The two-parameter determinant in both invariant coordinates and the
full three-variable Jacobian; the original-map normalization and collision;
diagonal equivalence, degrees, support, and the universal collision family.
\item The two degenerate axes, the two-sided inverse identity with a formal
arbitrary function $H$, the source-shear identities and sample degree
profiles, and the one-variable coefficient identities and explicit endpoint
of the classification proof.
\item The invertible variable-parameter scheme normalization, the eighteen
coefficient equations, the rank minor, the tangent
vector, the cokernel functional, the second-order obstruction, and the
vanishing of every coefficient equation under the dual-number presentation.
\end{enumerate}
The last check proves one ideal containment by direct substitution. The
reverse containment is justified in the article by the classification,
rank, and obstruction argument; the script does not silently label a
one-sided ideal check as equality.

The run included with this report used SymPy~1.14.0 and passed all thirteen
check groups. The JSON certificate includes the full rational matrix and
both row and column orderings. No numerical tolerance, random test, or
floating-point approximation is used. These checks are reproducibility
evidence for the finite algebra; they do not replace the universal
one-variable degree proofs.

\subsection{A realistic proof-assistant boundary}

The new theorems in this report have not been compiled in Lean or Rocq.
The existing repository formalizations of the original map are separate
from the new classification. A sensible formalization order is the
following dependency chain:
\[
\begin{gathered}
\text{weighted monomial encoding and localization identity}\\
\Downarrow\\
\text{constant weighted Wronskian lemma and cubic--square lemma}\\
\Downarrow\\
\text{three coefficient equations and exhaustive case split}\\
\Downarrow\\
\text{explicit classification, inverse, and collision theorems}\\
\Downarrow\\
\text{finite matrix certificates and double-point criterion}.
\end{gathered}
\]
The main classification needs polynomial degrees, characteristic-zero
derivatives, a rational-function constant lemma, and unique factorization
in one variable. It does not require the heavy geometric machinery used
in the literature to study unrestricted Keller maps. The double-point
statement additionally needs finite-type algebra, the Nullstellensatz,
Nakayama's lemma, and faithful scalar extension, or an alternative direct
ideal-membership certificate.

One attractive alternative for formalizing the last step is to produce
explicit polynomial combinations expressing the twelve simple generators
of the double-point ideal in terms of the eighteen original equations.
The present argument proves these combinations exist but does not supply
their expanded coefficients. Such certificates could isolate the formal
scheme result from a large commutative-algebra library.

\section{Further research questions and proposed next steps}
\label{wkr:sec:future}

The following questions are proposed by the present results. Unless
explicitly attributed, they are research directions formulated here, not
claims that a published source listed them as open. They are ordered to
separate immediately testable extensions from much broader classification
problems.

\subsection{Allow quadratic dependence on the second invariant}
\label{wkr:as:q:quadratic}% label added in batch 44

Keep $r$ affine but allow $\deg_vp,\deg_vq\le2$. Does the nonconstant
branch admit genuinely new diagonal-equivalence classes? After the change
$s=r$, the determinant becomes a polynomial of higher degree in $s$.
The highest coefficient should impose a weighted differential relation on
the leading coefficients, but the lower equations can couple several
common factors. The first concrete task is to derive the complete
coefficient system and identify whether an analogue of
Lemma~\ref{wkr:lem:wronskian} gives an a priori degree bound. A finite search
would be useful evidence, but an all-degree conclusion needs such a bound
or another structural argument.
[Added 29 September 2026, batch~44: still open. Part~II's
Proposition~\ref{wkr:as:prop:quadratic-repair} shows that a $v^2$ term is
needed, and suffices, to lift its obstructed slope directions at the tame
point to second order; that is not a classification of this class.]

\subsection{\texorpdfstring{Allow a nonconstant second-invariant coefficient in $r$}{Allow a nonconstant second-invariant coefficient in r}}
\label{wkr:as:q:slope}% label added in batch 44

Theorem~\ref{wkr:thm:triangular} settles arbitrary polynomial $t$-dependence
when the coefficient of $v$ in $r$ is constant. The next case is
$r=R(t)+G(t)v$ with nonconstant $G(t)$. Must $G$ be constant under suitable
Keller and polynomiality hypotheses, or do genuinely new normal forms
appear? The substitution $s=r$ now requires division by $G$ and is not a
global polynomial coordinate change. Track the behavior at the zeros of
$G$ and determine whether the constant-Wronskian argument has a localized
analogue. Allowing a quadratic term in $v$ is a further, distinct extension.
[Added 29 September 2026, batch~44: \emph{answered}. With $p,q$ affine in
$v$, $G$ must be constant, in all degrees
(Theorem~\ref{wkr:as:thm:slope}); the zeros of $G$ are controlled by a
rational pole equation (Lemma~\ref{wkr:as:lem:pole}). No new normal forms
appear (Theorem~\ref{wkr:as:thm:classification}). The quadratic extension
remains open (Section~\ref{wkr:as:q:quadratic}).]

\subsection{Does the nonreduced structure stabilize with degree?}
\label{wkr:as:q:stabilize}% label added in batch 44

For each $d\ge7$, form the normalized coefficient scheme with
$r=1+t+v$, $p,q$ affine in $v$, and lifted ordinary degree at most $d$.
Its reduced scheme is a point by Corollary~\ref{wkr:cor:reduced}. Is the full
scheme always the same double point? Or can high-degree coefficients
support new nilpotent directions invisible at degree seven? A first step
is to linearize the three identities in $(A,B,C,D)$ around the explicit
solution without truncating their degrees, solve the resulting polynomial
linear differential system, and then examine its obstruction maps.
This is a particularly well-targeted extension of the present work.
[Added 29 September 2026, batch~44: still open. Part~II's nilpotent
slopes (Theorem~\ref{wkr:as:thm:dual},
Corollary~\ref{wkr:as:cor:schemes}) live at the tame point with $r$ free,
not in this scheme with $r=1+t+v$ fixed; they neither extend nor
contradict Theorem~\ref{wkr:thm:double}.]

\subsection{Scheme-theoretic gluing at the invertible boundary}
\label{wkr:as:q:gluing}% label added in batch 44

Theorem~\ref{wkr:thm:torusdouble} settles the full degree-seven coefficient
scheme on $\beta\gamma\ne0$: it is the parameter torus times a double
point. Determine how this thickening meets the constant-$r$ tame locus
where that normalization ceases to be invertible. Does its scheme-theoretic
closure acquire embedded components or higher nilpotent order on either
axis or at their intersection? The reduced affine-plane closure and the
nodal boundary in Theorem~\ref{wkr:thm:plane} do not decide this gluing problem.
A concrete first step is to saturate the coefficient ideal by
$\beta\gamma$, then study its specialization on each axis with explicit
ideal-membership certificates.
[Added 29 September 2026, batch~44: still open; Part~II's
Theorem~\ref{wkr:as:thm:obstruction} constrains the tame end of the
gluing (Section~\ref{wkr:as:sec:future-glue}).]

\subsection{Other weights and other Wronskian exponents}

Our determinant produces a cubic--square relation because the first two
weights are $2$ and $1$. For other mixed-sign weight triples, analogous
highest-coefficient equations may have the form
$mCA'-nAC'=0$. Investigate the resulting common-power factorizations and
constant weighted Wronskians. The interesting issue is not just repeating
the calculation with new constants: the polynomiality inequalities at the
coordinate hyperplane may either permit or rule out the required
factorizations. A systematic analysis could delineate entire rigid classes
of equivariant maps.

\subsection{Compare lower-degree maps outside this symmetry class}

The current literature already includes a three-dimensional degree-four
example \cite{gao}; the existence of such an example is therefore not
posed as a new open problem here. (\wtag{} Gao's degree is the geometric
degree; the ordinary degrees of his map are $4,11,12$, see
Section~\ref{wkr:sub:gao}. A three-variable Keller map of \emph{ordinary}
degree below seven is not known from the sources cited here.) Instead, extract its invariant or
geometric structure and determine precisely which hypothesis of our
classification fails. Compare support size, rational coefficient height,
and the cost of a formal collision certificate. This would help replace a
single ``simplest map'' question by explicit, noninterchangeable complexity
criteria.

\subsection{Arithmetic optimization within the diagonal orbit}

Over $\Q$, the classification makes every nonconstant map in this class
explicitly diagonal-equivalent. One can now optimize primitive coefficient
height, denominator size, or integral collision height within that orbit.
The parameters $a,b$ and the additional unnormalizing target scalings make
these concrete arithmetic optimization problems. Prove lower bounds for a
specified normalization rather than relying on finite samples of small
rational parameters.

\subsection{Direct ideal certificates and kernel verification}

Generate explicit membership certificates for the twelve simple generators
in Appendix~\ref{wkr:app:presentation}, and check them with polynomial
normalization in Lean or Rocq. This would turn the deformation theorem into
a compact finite certificate after the all-degree classification has been
formalized. The target should be a transparent theorem boundary: no hidden
use of a computer-algebra claim that a Gr\"obner basis is complete.

\subsection{Positive characteristic as a controlled contrast}

The displayed family and finite identities have meaningful reductions at
many primes, but the all-degree proof does not automatically survive.
In positive characteristic, nonconstant $p$th powers have derivative zero,
and the weighted degree coefficients can vanish. Classify which parts of
the theorem persist under explicit degree bounds and which fail through
Frobenius phenomena. Separately determine the primes at which the tangent
rank or obstruction changes. Reducing a characteristic-zero certificate
is not sufficient to claim the same unrestricted classification.

\section{Conclusion}

The repository's finite sparsity problem admits a stronger structural
answer in its affine-invariant class. Constant third numerator gives all
and only the explicit tame maps \eqref{wkr:eq:tamepq}. Nonconstant third
numerator forces both invariant coefficients to be nonzero and forces the
entire map into the diagonal orbit \eqref{wkr:eq:family}, with degree seven and
sixteen terms. An explicit source shear extends the classification to
$r=R(t)+\gamma v$ and gives the exact degree spectrum
$\{7,8,10,12,\ldots\}$. The proof is an all-degree argument, not an
extrapolation from a coefficient search.

The parameter-space results explain two different notions of rigidity.
Before fixing the third numerator, the nonconstant branch has an affine
plane closure with invertible coordinate axes. After fixing $r=1+t+v$,
all reduced families are constant, yet the degree-seven coefficient scheme
retains exactly one square-zero parameter. Its tangent direction has an
explicit nonzero second-order obstruction. Restoring the two nonconstant
parameters gives the torus times this double point. These distinctions are essential
both for further searches and for a trustworthy formalization.

% ==== Part II begins
%% ====================================================================
%% Batch-44 addition (manuscript 04, "All-Degree Classification of
%% Affine-in-z Weighted Keller Maps"). All labels of Part II carry the
%% prefix wkr:as: (affine slope). Manuscript section k is printed here as
%% Section k+12; its Appendices A and B are Sections 24 and 25.
%% ====================================================================
\clearpage
\phantomsection\part*{Part II.\quad All-degree classification of affine-in-\texorpdfstring{$z$}{z} weighted Keller maps}
\addcontentsline{toc}{part}{Part II.\texorpdfstring{\quad}{ } Affine-in-\texorpdfstring{$z$}{z} weighted Keller maps}
\fancyhead[L]{\small\color{inkblue}Part II: affine-in-$z$ weighted Keller maps}

\noindent\emph{Sections~1--11 and Appendices~A--D form
the original report (Part~I) of 24 September 2026, unchanged except for
dated pointers to this addition. Part~II was added on 29 September 2026 in
batch~44 of ProveIt's incoming-report intake, from one manuscript written
against Part~I as printed above. It answers the question of
Section~\ref{wkr:as:q:slope} affirmatively: when all three invariant
numerators are affine in $v$, the Keller condition forces the coefficient
of $v$ in $r$ to be constant, with no bound on ordinary degree.
Section~\ref{wkr:as:sec:provenance} records the provenance, fixes the
notation against Part~I and states what is answered and what is not.}

\section{Provenance, conventions and what Part~II answers}
\label{wkr:as:sec:provenance}

\subsection{One manuscript, one addition}\label{wkr:as:sec:source}
Part~II is batch-44 manuscript~04, dated 29~September 2026:
\emph{All-Degree Classification of Affine-in-$z$ Weighted Keller Maps:
Constant-slope rigidity and obstructed nilpotent deformations}. Its author
line is ``Prepared for Vladimir Reshetnikov'' (PDF metadata:
``AI-assisted research draft prepared for Vladimir Reshetnikov''). It was
delivered as the archive
\texttt{ProveIt\_Affine\_Weighted\_Keller\_Classification} (inner directory
of the same name, main file \texttt{article.tex}, a 22-page A4 PDF). It is
pinned to the ProveIt commit
\begin{center}
\texttt{9b24a3a8d545af9624f6ac455f5b548be62818b6}
\end{center}
(committed 2026-09-29 16:02:05~UTC). The archive arrived in commit
\texttt{ae28ea2db} and was placed in \texttt{203016015} as an
\emph{addition} to this report rather than as a new report, because it
answers Part~I's Section~\ref{wkr:as:q:slope} and restates Part~I's normal
forms. Neither Part~I nor the Jacobian project
\rpath{Algebra/JacobianConjecture} changed between the pin and the
placement commit, so the manuscript read exactly the text printed above.

Its companion files are shipped with the prefix \texttt{02-affine-slope-}
(Section~\ref{wkr:as:sec:repro}); its manuscript, README and PDF are not
shipped, this Part being their written form. Every result, proof, example,
table, question and limitation of the manuscript is printed; the few items
that duplicate Part~I are printed once
(Section~\ref{wkr:as:sec:duplicates}). Text added at the write, beyond the
present section, is marked \textup{[write]}. The manuscript's title page
and abstract read as follows.

\begin{quote}\small
\textbf{Main result.} In characteristic zero, the Keller condition forces
the coefficient of $v=x^2z$ in the third invariant numerator to be
constant, provided all three invariant numerators are affine in $v$. This
closes the explicit extension question in Section~10.2 of ProveIt's
\emph{All-Degree Rigidity of a Weighted Keller Class} (Part~I). The
resulting classification has no ordinary degree bound.

\textbf{A second result.} The field-level rigidity does not hold over dual
numbers. Nonconstant nilpotent slopes exist in arbitrarily high degrees,
but every nonzero slope direction at the simplest tame map is obstructed
at second order, even when arbitrary affine-in-$v$ second-order
corrections are allowed.

\textbf{Status.} Complete written proofs and reproducible exact symbolic
checks are supplied. The new results are not formalized in Lean or Rocq,
not independently refereed, and not asserted to have established
worldwide priority. The known counterexample and the previously
classified constant-slope family are credited separately.

\textbf{Abstract.} Let $K$ be a field of characteristic zero, let $t=xy$
and $v=x^2z$, and consider polynomial lifts
\[
F(x,y,z)=\left(x^{-2}p(t,v),\,x^{-1}q(t,v),\,xr(t,v)\right),
\qquad \deg_v p,\deg_v q,\deg_v r\leq1.
\]
These are exactly the maps affine in $z$ with source weights $(-1,1,2)$
and target weights $(2,1,-1)$. We prove that a nonzero constant Jacobian
forces $r_v$ to be constant. The proof is not a bounded coefficient
search: a highest-coefficient identity gives a cube--square
factorization; a rational moving coordinate produces an equation
$LL'=\alpha/U^3$; pole orders and polynomiality at the coordinate
hyperplane then exclude every nonconstant slope. Combined with an
explicit reconstruction, this gives all maps in the class: a tame family,
or one fixed noninjective map up to diagonal scalings and a single
equivariant source shear. We supply an exact coefficient-only recognition
algorithm, inverses or collision certificates, and the degree spectrum
$\{7\}\cup\{8,10,12,\ldots\}$ for the noninjective branch.

We also determine a sharp limitation of this rigidity. It persists over
reduced $\Q$-algebras, but fails over dual numbers. At the map $(z,y,x)$,
any nonzero first-order variation of the third slope has a nonvanishing
second-order obstruction, independent of all second-order corrections
inside the affine-in-$z$ class. In bounded coefficient schemes, this gives
nonzero nilpotent slope coordinates and arbitrarily large spaces of
obstructed tangent directions. Full proofs, source comparisons, twelve
exact-check groups, and a further research program are included.
\end{quote}

\subsection{What Part~II answers, and what it does not}
\label{wkr:as:sec:answers}
\begin{itemize}
\item \emph{Section~\ref{wkr:as:q:slope} is answered, affirmatively.}
With $p,q$ affine in $v$ and $r=R(t)+G(t)v$ in Part~I's notation,
Theorem~\ref{wkr:as:thm:slope} proves that $G$ is constant, with no bound
on ordinary degree. Part~I's Theorem~\ref{wkr:thm:triangular} then
applies, so Part~I's classification covers every affine-in-$z$ weighted
Keller map (Theorem~\ref{wkr:as:thm:classification}). The exclusions
stated at the end of Section~\ref{wkr:sec:triangular} and in
Section~\ref{wkr:sec:verification} are lifted accordingly (dated pointers
there).
\item \emph{Section~\ref{wkr:as:q:quadratic} ($v^2$ terms in $p,q$) is not
answered.} Proposition~\ref{wkr:as:prop:quadratic-repair} shows that a
$v^2$ term is needed, and suffices, to lift the canonical obstructed
directions at the tame point to second order; it is not a classification
of the quadratic class.
\item \emph{Section~\ref{wkr:as:q:stabilize} is not answered.} The
nilpotent results of Section~\ref{wkr:as:sec:deformations} concern slope
directions at the tame point $(v,t,1)$, in coefficient schemes in which
$r$ varies; Theorem~\ref{wkr:thm:double} concerns the noninjective point
with $r=1+t+v$ fixed. They neither extend nor contradict the double point,
as the manuscript says itself (Section~\ref{wkr:as:sec:future-thick}).
\item \emph{Section~\ref{wkr:as:q:gluing} is not answered};
Theorem~\ref{wkr:as:thm:obstruction} is a constraint on it
(Section~\ref{wkr:as:sec:future-glue}).
\end{itemize}

\subsection{Notation: Part~I against Part~II}\label{wkr:as:sec:notation}
Part~II keeps the manuscript's letters, which were chosen before it was
joined to Part~I, so several letters mean different things in the two
Parts. Table~\ref{wkr:as:tab:notation} fixes them; within Part~II only the
Part~II column applies. The tempting false readings are printed in the
last column.

\begingroup\small
\setlength{\tabcolsep}{4pt}
\begin{longtable}{@{}>{\raggedright}p{0.10\textwidth}>{\raggedright}p{0.27\textwidth}>{\raggedright}p{0.31\textwidth}>{\raggedright\arraybackslash}p{0.24\textwidth}@{}}
\caption{Symbols whose meaning differs between the Parts.}
\label{wkr:as:tab:notation}\\
\toprule
Symbol & Part I & Part II & Not to be read as\\
\midrule
\endfirsthead
\toprule
Symbol & Part I & Part II & Not to be read as\\
\midrule
\endhead
$a,b$ & family parameters of $F_{a,b}$, $a=\beta^2/\gamma$, $b=\gamma/\beta$ &
coefficient polynomials, $p=a(t)v+b(t)$; Part~I's parameters are written
$a_{\mathrm I},b_{\mathrm I}$ & the parameters of \eqref{wkr:eq:family}\\
$c,d$ & scalars in $C=cE^2$, $D=dt-c$ \eqref{wkr:eq:linearforms}; $c$ also
the coefficient vector, $\cstar$ its point & coefficient polynomials,
$q=c(t)v+d(t)$ & scalars\\
$e,g$ & $e$: the square-zero element of Corollary~\ref{wkr:cor:functor} &
$r=g(t)v+e(t)$: $e=r(t,0)$, and $g=r_v$ is the \emph{third slope}; they
are Part~I's $R(t)$ and $G(t)$ of Section~\ref{wkr:as:q:slope} & nilpotents\\
$G$ & coefficient of $v$ in $r$ (Sections~\ref{wkr:sec:triangular},
\ref{wkr:as:q:slope}); $G_{\mathrm{Gao}}$ is Gao's map & a first-order
slope variation (Theorems~\ref{wkr:as:thm:dual},
\ref{wkr:as:thm:obstruction}); the manuscript's valuation
$G=\ord(g)$ is written $o_g$; Section~\ref{wkr:as:sec:intro}
quotes Part~I's $G$, which is $g$ here & Gao's map\\
$A,B,C,D$ & $p=As+B$, $q=Cs+D$ \eqref{wkr:eq:ABCD} & the same in
\eqref{wkr:as:eq:moving}--\eqref{wkr:as:eq:movingcoeff}, with rational
$A,\ldots,D$; first-order variations of $a,b,c,d$ in
Theorem~\ref{wkr:as:thm:obstruction}. The manuscript's $A=\ord(a)$,
$C=\ord(c)$, its collision points $A,B$, its reduced algebra $A$ and its
coefficient algebras $A_D$ are written $o_a,o_c$, $X_1,X_2$, $\Lambda$ and
$\Lambda_{\Dcap}$ & the same objects in every section\\
$U$ & an output coordinate \eqref{wkr:eq:tameinverse} & a polynomial in
$t$; at the end $U=1-2\beta t/3$, which is Part~I's $E$ of
\eqref{wkr:eq:compactendpoint} & a coordinate\\
$k$ & $E=1+kt$ & the scalar in $c=kHU^2$; at the end $k=9/\beta$, which is
Part~I's scalar $c$ of \eqref{wkr:eq:endpoint} & Part~I's $k=-2\beta/3$\\
$H$ & $B=t^2H(t)$ \eqref{wkr:eq:tameF}; the residual vector
\eqref{wkr:eq:H} & the cube root in $g=H^3$ & either\\
$s,S$ & $s=1+\beta t+v$ & $s=e+H^3v$ and $S=e+\gamma v$; each is $r$ & a
constant\\
$L$, $\mathcal L$ & $L$: the row functional \eqref{wkr:eq:ell};
$\mathcal L$: the lift & $L=\delta-\alpha T$, a rational function;
$\mathcal L$ is the same lift & a functional\\
$\alpha$ & (write notes) the project's coefficient $[t^2]p$ & $\alpha=2k/3$
in \eqref{wkr:as:eq:VT} & $[t^2]p$\\
$T$ & only in $T_h$ & $T=B/w^2$ (Section~\ref{wkr:as:sec:rigidity}); the
indeterminate of the cubic \eqref{wkr:as:eq:cubic}; $T_h$ as in Part~I &
one object\\
$P,Q,R$ & $R(t)$: the $t$-part of $r$ in Theorem~\ref{wkr:thm:triangular}
& the components of $F$, and target coordinates in the cubic model;
Section~\ref{wkr:as:sec:intro} quotes Part~I's $R(t)$ &
Part~I's $R(t)$, which is $e$ here\\
$\Dcap$ & --- & the degree cap of $\Lambda_{\Dcap}$ and the input degree of
the algorithm (the manuscript writes $D$) & the polynomial $D$\\
$\tau$ & --- & $y+1/x$ (Section~\ref{wkr:as:sec:cubic}) & $t=xy$\\
\bottomrule
\end{longtable}
\endgroup

The field, the weights, the invariants $t=xy$, $v=x^2z$, the lift
$\mathcal L$, the Jacobian $\J$, the normalization to determinant $-1$,
the pair $\beta,\gamma$ (for $e=1+\beta t$ and $g=\gamma$), the source
shear $T_h$ and its polynomial $h$, the map $F_0$ and its normalization
$F_*$, and the dual-number variable $\eps$ mean the same in both Parts.
``Degree'' means ordinary polynomial degree unless qualified as
function-field or geometric degree. Renamed from the manuscript: its field
$\mathbb K$ is printed $K$; its valuations $A,C,G$ are $o_a,o_c,o_g$; its
collision points $A,B$ are $X_1,X_2$; its reduced algebra $A$ is
$\Lambda$; its degree cap $D$ and coefficient algebras $A_D$ are $\Dcap$
and $\Lambda_{\Dcap}$; Part~I's parameters in its comparison sentence are
$a_{\mathrm I},b_{\mathrm I}$. No normalization was changed. Equations are
numbered by section, as in Part~I; the manuscript numbered them
consecutively, (1)--(59), and its section~$k$ is Section~$k+12$ here.

\subsection{What is printed once}\label{wkr:as:sec:duplicates}
The manuscript restates several items of Part~I to be self-contained. They
are printed once, and where they occur Part~II points to them.
\begin{itemize}
\item The map $F_0$, its determinant $-2$ and its integral collision
(manuscript equations (1)--(2)) are \eqref{wkr:eq:original},
\eqref{wkr:eq:oldcollision}; the normalization $F_*=\diag\circ F_0\circ
\diag$ (manuscript equation (39)) is \eqref{wkr:eq:originalnormalization}.
\item The lift and the weighted Jacobian identity with its chain-rule
proof (manuscript equation (4) and Proposition~2.2 with equation (7)) are
\eqref{wkr:eq:lift} and Proposition~\ref{wkr:prop:jacobian}, which are in
turn the project's (Section~\ref{wkr:sub:project}). The manuscript's
coefficient expansion \eqref{wkr:as:eq:J0}--\eqref{wkr:as:eq:J2} is new
and printed.
\item The tame inverse and its decomposition into two elementary shears
and a permutation (manuscript Section~6.1, equations (34)--(35), with output
coordinates $(X,Y,Z)$) are \eqref{wkr:eq:tameF}--\eqref{wkr:eq:tameinverse}
(output coordinates $(U,V,W)$).
\item The proof of the degree spectrum (manuscript Corollary~6.4) repeats
Part~I's proof of Corollary~\ref{wkr:cor:degreespectrum}; the statement is
printed as Corollary~\ref{wkr:as:cor:degrees} because its scope is larger,
and its proof refers to Part~I.
\item Printed although partly known, because they carry the argument:
the determinant expansion \eqref{wkr:as:eq:movingcoeff}, which is Part~I's
\eqref{wkr:eq:threeidentities} with \eqref{wkr:eq:E0}--\eqref{wkr:eq:E2}
(here with rational coefficients); Lemma~\ref{wkr:as:lem:constants},
which Part~I proves inside Lemma~\ref{wkr:lem:powers}; the normal forms
of Theorem~\ref{wkr:as:thm:classification}, which restate
Theorems~\ref{wkr:thm:classification} and~\ref{wkr:thm:triangular} in the
coordinates $\beta,\gamma,e$ over the enlarged class
(Remark~\ref{wkr:as:rem:partI}); and the cubic model of
Section~\ref{wkr:as:sec:cubic}, which the manuscript does not advertise
as new (see the \textup{[write]} note there).
\end{itemize}
The manuscript's reconstruction of the noninjective branch
(Section~\ref{wkr:as:sec:reconstruction}) is a genuinely different route
from Part~I's Section~\ref{wkr:sec:proof}: it reuses the moving-coordinate
equations with $H=1$ instead of the constant weighted Wronskian lemmas.
It is printed in full as a second proof.

\subsection{The formal project and Part~II}\label{wkr:as:sec:formal}
Check group~4 of the companion script (\texttt{original\_map}) rechecks in
SymPy statements that the project has \emph{kernel-checked} (Lean
namespace \texttt{LeanProofs.JacobianCounterexample}; verified at the
write in the files at the placement commit):
\begin{itemize}
\item $\det DF_0=-2$: Lean \texttt{jacobianDet\_counterexample}
(\rpath{Algebra/JacobianConjecture/Lean/JacobianConjecture/Counterexample.lean},
line~134), a polynomial identity over every commutative ring; Rocq
\texttt{jacobian\_det\_is\_minus\_two}
(\rpath{Algebra/JacobianConjecture/Coq/Counterexample.v}, line~182), at
every point;
\item the collision $F_0(-1,1,5)=F_0(0,-2,-16)=(0,-2,0)$: Lean
\texttt{collision\textsubscript{0}}, \texttt{collision\textsubscript{1}},
\texttt{collisionValue}, \texttt{collision} and
\texttt{collision\_points\_distinct} (\rpath{Counterexample.lean}, lines
143--170); Rocq \texttt{integral\_collision\_0\_value},
\texttt{integral\_collision\_1\_value} and
\texttt{integral\_collision\_points\_distinct} (\rpath{Counterexample.v},
lines 249--265); and the refutation
\texttt{jacobianConjectureInDimensionThree\_false}
(\rpath{Counterexample.lean}, line~201; Rocq
\texttt{jacobian\_conjecture\_dimension\_three\_is\_false},
\rpath{Counterexample.v}, line~331).
\end{itemize}
The weighted torus action on which the whole class rests is
\texttt{counterexample\_scaling}
(\rpath{Algebra/JacobianConjecture/Lean/JacobianConjecture/Scaling.lean},
line~35; Rocq \rpath{Scaling.v}, line~32). The same group also checks the
normalization \eqref{wkr:eq:originalnormalization}, which is \emph{not}
formalized, and nothing else of Part~II is: not the weighted Jacobian
identity, not the normal forms (T) and (N), and none of the theorems,
propositions or corollaries of Sections~\ref{wkr:as:sec:main}--%
\ref{wkr:as:sec:algorithm}. The package contains no Lean or Rocq file.
Placement of this report beside the Lean/Rocq development confers no
formal status on it; the manuscript says so correctly (its source ledger:
the package ``neither edits the repository nor inherits the formal status
of declarations stored elsewhere in it'').
\section{The question, the contribution, and its boundaries}
\label{wkr:as:sec:intro}

\subsection{A specific unresolved question in the repository}

ProveIt's Jacobian project records the polynomial map $F_0$ of
\eqref{wkr:eq:original}. The project attributes the map to Alp\"oge's July
2026 announcement and contains Lean and Rocq developments for its constant
Jacobian and explicit collisions \cite{repo-9b24}. The identities needed
here are elementary: $\det DF_0=-2$ and the collision
\eqref{wkr:eq:oldcollision}. We check them independently in the companion
code. They are not new results of this article.
\textup{[write]} The map and these two identities are printed once, in
Part~I; their kernel-checked status is in
Section~\ref{wkr:as:sec:formal}.

A subsequent repository report, \emph{All-Degree Rigidity of a Weighted
Keller Class} (Part~I), classifies a normalized weighted class when the
third invariant numerator is $r=R(t)+\gamma v$ with $\gamma\in K$. Its
Section~\ref{wkr:as:q:slope} explicitly asks whether allowing
\[
r=R(t)+G(t)v
\]
with a nonconstant polynomial $G$ produces new normal forms, or whether
the Keller and polynomiality conditions force $G$ to be constant. The
obstacle identified there is real: the substitution $s=r$ divides by $G$
and is not a global polynomial coordinate change.

Theorem~\ref{wkr:as:thm:slope} resolves that question under the report's
retained hypotheses on $p$ and $q$: \emph{$G$ must be constant}. No bound
on the ordinary degree is added. The proof explicitly controls the zeros
of $G$; it never treats the rational substitution as a polynomial
automorphism.

\subsection{What is new here, and what is inherited}

The new field-level result is constant-slope rigidity for the full
three-numerator affine-in-$v$ class. The new infinitesimal results are the
reduced-base extension, nonconstant dual-number examples, universal
second-order obstruction at the simplest tame point, and the resulting
nilpotent directions in finite coefficient schemes.

The constant-slope normal forms, their tame inverse, diagonal
equivalence, source-shear extension, core support profile, and degree
spectrum already occur in Part~I. We rederive them after proving the
missing rigidity step, both to make this article self-contained and to
extend their \emph{scope} to every affine-in-$z$ map with the specified
weights. The weighted determinant identity and its localization proof were
already present in the project's research notes, as acknowledged by that
report. The cubic function-field calculation in
Section~\ref{wkr:as:sec:cubic} is likewise not advertised as new for
$F_0$; its role is to make the geometric consequences of the enlarged
classification explicit.

\subsection{Relation to the current primary literature}

Shaska's revised paper \cite{shaska} studies the graded structure of the
counterexample, quotient coordinates, and coefficient schemes. Gao
\cite{gao} develops a broader construction of noninjective Keller maps
from curves, including maps of different geometric degrees. Those works
supply important context for why a structural classification inside a
specified graded class is useful; neither is needed as a black-box
classification theorem in our proof.

The source check for this article used the repository pin
(Section~\ref{wkr:as:sec:source}), the precise prior report and its stated
exclusions, Shaska's version~2, and Gao's version~1. This is a targeted
comparison, not an exhaustive bibliographic priority determination. In
particular, the present work is \emph{not} a new disproof of the Jacobian
conjecture, a classification of arbitrary Keller maps, or a
minimum-degree theorem without symmetry and linearity hypotheses.
``Degree'' below always means ordinary polynomial degree unless
explicitly qualified as function-field or geometric degree.

\section{Weighted coordinates and the exact coefficient equations}
\label{wkr:as:sec:coordinates}

Throughout Sections~\ref{wkr:as:sec:coordinates}--\ref{wkr:as:sec:algorithm},
$K$ is a field of characteristic zero. We write primes for
differentiation with respect to $t$.

\begin{definition}
A map $F=(P,Q,R):\A^3_K\to\A^3_K$ is \emph{weighted} in this article if
\begin{equation}\label{wkr:as:eq:weights}
F(\rho^{-1}x,\rho y,\rho^2z)
 =\bigl(\rho^2P(x,y,z),\rho Q(x,y,z),\rho^{-1}R(x,y,z)\bigr)
\end{equation}
as an identity with an invertible indeterminate $\rho$. It is
\emph{Keller} if $\det DF\in K^\times$.
\end{definition}

Set $t=xy$ and $v=x^2z$. A monomial $x^iy^jz^k$ of weight $2$ satisfies
$i=j+2k-2$, so multiplying it by $x^2$ gives $t^jv^k$. The corresponding
exponents for weights $1$ and $-1$ are $j+2k-1$ and $j+2k+1$.
Consequently every map satisfying \eqref{wkr:as:eq:weights} has a unique
expression $F=\mathcal L(p,q,r)$ in the localization at $x$, as in
\eqref{wkr:eq:lift}. It is affine in $z$ in each coordinate if and only if
\begin{equation}\label{wkr:as:eq:abcdef}
p=a(t)v+b(t),\qquad q=c(t)v+d(t),\qquad r=g(t)v+e(t)
\end{equation}
with all six coefficients in $K[t]$. The lift is polynomial exactly when
\begin{equation}\label{wkr:as:eq:polynomiality}
b\in t^2K[t],\qquad d\in tK[t].
\end{equation}
Indeed, a monomial $t^iv^j$ in $p$ or $q$ must satisfy respectively
$i+2j\ge2$ or $i+2j\ge1$; for $j\le1$ these are precisely
\eqref{wkr:as:eq:polynomiality}. There is no analogous restriction on $e$
or $g$.

For every polynomial lift, $\det DF=\J(p,q,r)$ with $\J$ the weighted
determinant \eqref{wkr:eq:J} evaluated at $(t,v)=(xy,x^2z)$
(Proposition~\ref{wkr:prop:jacobian}).
\textup{[write]} The manuscript states and proves this again as its
Proposition~2.2, by the same chain-rule computation in the localization
at $x$; it is printed once, in Part~I, and is the project's
(Section~\ref{wkr:sub:project}).

The substitution $K[t,v]\to K[x,y,z]$ is injective: distinct $t^iv^j$
produce distinct monomials $x^{i+2j}y^iz^j$. Thus we may work with $\J$
as a polynomial in $t,v$. Its value at $t=v=0$ is
\begin{equation}\label{wkr:as:eq:normalization}
\J(0,0)=-a(0)d'(0)e(0).
\end{equation}
For a Keller map, all three factors are nonzero. Dividing the first,
second, and third target coordinates respectively by these factors
normalizes
\begin{equation}\label{wkr:as:eq:normalized}
a(0)=d'(0)=e(0)=1,\qquad \J=-1.
\end{equation}
This normalization loses no Keller map and preserves the structural class.

For reference, direct expansion gives
\begin{equation}\label{wkr:as:eq:Jcoeff}
\J=J_0+J_1v+J_2v^2,
\end{equation}
where
\begin{align}
J_0={}&(ce+dg)b'+2b(ce'-gd')-a(de)',\label{wkr:as:eq:J0}\\
J_1={}&(ce+dg)a'+2b(cg'-gc')\notag\\
 &\quad+a(ce'-dg'-ec'-3gd')+2cgb',\label{wkr:as:eq:J1}\\
J_2={}&2cga'-3agc'+acg'.\label{wkr:as:eq:J2}
\end{align}
These identities will also make the deformation calculation transparent.

\section{The main classification}
\label{wkr:as:sec:main}

\begin{theorem}[Constant-slope rigidity]\label{wkr:as:thm:slope}
Let \eqref{wkr:as:eq:abcdef} satisfy polynomiality
\eqref{wkr:as:eq:polynomiality} and the normalized Keller identities
\eqref{wkr:as:eq:normalized}. Then $g$ is constant. Equivalently, every
weighted Keller map affine in $z$ has a constant third invariant slope
$r_v$, after or before target normalization.
\end{theorem}

\begin{theorem}[Complete normal forms]\label{wkr:as:thm:classification}
Under the hypotheses of Theorem~\ref{wkr:as:thm:slope}, exactly one of the
following two cases occurs.

\smallskip
\noindent\textup{(T) Tame branch.} There are unique
$\lambda\in K$ and $B\in t^2K[t]$ such that
\begin{equation}\label{wkr:as:eq:tame-normal}
p=v+B(t),\qquad q=t+\lambda\bigl(v+B(t)\bigr),\qquad r=1.
\end{equation}

\smallskip
\noindent\textup{(N) Noninjective branch.} There are unique
$\beta,\gamma\in K^\times$ and $e\in K[t]$ with
$e(0)=1$ and $e'(0)=\beta$ such that, on setting
\begin{equation}\label{wkr:as:eq:US}
U=1-\frac{2\beta}{3}t,\qquad S=e(t)+\gamma v,
\end{equation}
one has
\begin{align}
p&=\frac{U^3S-(U^2+U)/2}{\gamma},\label{wkr:as:eq:normalp}\\
q&=\frac9\beta\left(U^2S-\frac{2U+1}{3}\right),\label{wkr:as:eq:normalq}\\
r&=S.\label{wkr:as:eq:normalr}
\end{align}
Conversely every displayed triple gives a normalized polynomial Keller
map. The maps in \textup{(T)} are tame automorphisms, and those in
\textup{(N)} are noninjective over the ground field $K$.
\end{theorem}

Here no ordinary degree bound is imposed on $a,b,c,d,e,g$. The
classification is therefore not an extrapolation from finitely many
coefficient schemes. The parameters of (N) are read directly from $r$:
$\gamma=r_v$, $e=r(t,0)$, and $\beta=e'(0)$. The inverse and collision
assertions are proved in Section~\ref{wkr:as:sec:consequences}.

\begin{remark}[\textup{[write]} Relation to Part~I]\label{wkr:as:rem:partI}
Branch (T) is Part~I's case~\ref{wkr:case:tame}, equation
\eqref{wkr:eq:tamepq}, verbatim. For $e=1+\beta t$, branch (N) is Part~I's
$F_{a_{\mathrm I},b_{\mathrm I}}$ of \eqref{wkr:eq:family} with
$a_{\mathrm I}=\beta^2/\gamma$, $b_{\mathrm I}=\gamma/\beta$: $U$ is
Part~I's $E$ and \eqref{wkr:as:eq:normalp}--\eqref{wkr:as:eq:normalq}
expand to \eqref{wkr:eq:compactendpoint} before the $\gamma$-rescaling,
since $-(U^2+U)/2=-1+\beta t-\tfrac{2\beta^2}{9}t^2$ and
$-\tfrac9\beta\cdot\tfrac{2U+1}{3}=4t-\tfrac9\beta$. For general $e$ it is
Part~I's $F_{a_{\mathrm I},b_{\mathrm I}}\circ T_h$ of
Theorem~\ref{wkr:thm:triangular} (see \eqref{wkr:as:eq:core-shear}). What
is new is that Theorem~\ref{wkr:as:thm:slope} makes these the \emph{only}
forms once $r$ is merely affine in $v$. The normal form (N) with symbolic
$\beta,\gamma,e$ was rechecked at the write (Jacobian $-1$ identically).
\end{remark}

\section{Proof of constant-slope rigidity}
\label{wkr:as:sec:rigidity}

The proof separates three issues that a formal coefficient search would
not settle: arbitrary degree cancellations, poles of the moving
coordinate, and polynomiality across the coordinate hyperplane.

\subsection{Two elementary rational-function facts}

\begin{lemma}[Constants of differentiation]\label{wkr:as:lem:constants}
If $f\in K(t)$ and $f'=0$, then $f\in K$.
\end{lemma}
\begin{proof}
Write $f=A/B$ with coprime polynomials. The identity $A'B-AB'=0$ implies
$A\mid A'$ if $A\ne0$ and $B\mid B'$. A nonconstant polynomial in
characteristic zero has derivative of strictly smaller degree, so both
must be constant. The zero function is immediate.
\end{proof}

\begin{lemma}[A rational pole equation]\label{wkr:as:lem:pole}
Let $K$ be algebraically closed of characteristic zero. Suppose
$U\in K[t]\setminus\{0\}$, $L\in K(t)$, and
$\alpha\in K^\times$ satisfy
\begin{equation}\label{wkr:as:eq:pole}
LL'=\frac{\alpha}{U^3}.
\end{equation}
Then for some $\xi\in K$ and nonzero scalars $u,\ell$,
\[
U=u(t-\xi)^m,\qquad L=\ell(t-\xi)^{-n},
\qquad m,n\ge1,\qquad 3m=2n+1.
\]
If, in addition, $H\in K[t]$ is coprime to $U$ and
$\frac UH(\delta_0-L)$ is regular at $\xi$ for a scalar $\delta_0$,
then $m=n=1$.
\end{lemma}
\begin{proof}
The right side is nonzero, so $L$ and $L'$ are nonzero. If $L$ had a
finite zero of order $j\ge1$, the order of $LL'$ there would be $2j-1>0$.
The right side of \eqref{wkr:as:eq:pole} has order either $0$ or a
negative multiple of $3$. This is impossible. A rational function over an
algebraically closed field with no finite zero has constant numerator in
reduced form. Thus $L=\ell_0/V$ for a nonzero scalar $\ell_0$ and a
nonconstant polynomial $V$.

Substitution gives
\begin{equation}\label{wkr:as:eq:critical}
U^3V'=-\frac{\alpha}{\ell_0^2}V^3.
\end{equation}
Every root of $V'$ is therefore a root of $V$. If $V$ has degree $N$ and
$r$ distinct roots, the total multiplicity of those roots in $V'$ is
$N-r$: a root of multiplicity $j$ contributes exactly $j-1$ in
characteristic zero. Since all $N-1$ roots of $V'$ occur there, $r=1$.
Write $V=c(t-\xi)^n$. Equation~\eqref{wkr:as:eq:critical} then forces
$U=u(t-\xi)^m$ and $3m=2n+1$. In particular $m,n\ge1$.

Since $H(\xi)\ne0$ and $L$ has a pole, the order at $\xi$ of
$(U/H)(\delta_0-L)$ is exactly $m-n=(1-m)/2$. Regularity forces $m\le1$,
so $m=n=1$.
\end{proof}
\noindent\textup{[write]} Inside this proof only, $r$ counts distinct
roots and $c$ is a scalar; neither is the numerator $r$ or the
coefficient polynomial $c$.

\subsection{The nonzero-slope case has nonzero second slope}

Assume $g\ne0$. If $c=0$, the equation $J_1=0$ reduces to
\[
dga'-adg'-3agd'=0.
\]
Here $a,d,g$ are nonzero: $a(0)=1$ and $d'(0)=1$. Dividing in $K(t)$
shows that
\[
\left(\frac{a}{gd^3}\right)'=0.
\]
Lemma~\ref{wkr:as:lem:constants} gives $a=\kappa gd^3$ with
$\kappa\in K^\times$. But $d(0)=0$, contradicting $a(0)=1$.
Hence
\begin{equation}\label{wkr:as:eq:nonzeroacg}
a\ne0,\qquad c\ne0,\qquad g\ne0.
\end{equation}

\subsection{Cube--square factorization with coprime factors}

We may extend scalars to an algebraic closure for this part of the proof.
A polynomial proved constant after this extension was already constant
over the original field. Equation $J_2=0$ gives
\[
\left(\frac{a^2g}{c^3}\right)'=0,
\qquad \frac{a^2g}{c^3}\in K^\times
\]
by Lemma~\ref{wkr:as:lem:constants}. At a finite point, write
$o_a=\ord(a)$, $o_c=\ord(c)$, $o_g=\ord(g)$. Then
\begin{equation}\label{wkr:as:eq:valuationrelation}
2o_a+o_g=3o_c.
\end{equation}
The three polynomials $a,c,g$ cannot vanish at the same point: otherwise
the last column of the determinant \eqref{wkr:eq:J} would vanish there
for every $v$, contradicting $\J=-1$.

If $o_g>0$, equation~\eqref{wkr:as:eq:valuationrelation} and the
preceding observation force $o_a=0$ and $o_g=3o_c$. If $o_a>0$, then
$o_g=0$ and $2o_a=3o_c$, so $o_a$ is a multiple of $3$. Unique
factorization, including nonzero scalar factors, therefore gives
\begin{equation}\label{wkr:as:eq:UHfactor}
a=U^3,\qquad c=kHU^2,\qquad g=H^3,
\qquad (U,H)=1,
\end{equation}
for polynomials $U,H$ and a scalar $k\ne0$. We choose the cube root in $U$
so that $U(0)=1$, which is possible because $a(0)=1$. Choosing a cube
root of the scalar in $g$ is harmless over the algebraic closure. Scalar
descent will be automatic once $g$ is proved constant.

\subsection{A rational moving coordinate and its differential invariant}

Work temporarily over $K(t)$, and set
\[
w=\frac UH,\qquad s=e+H^3v.
\]
Then uniquely
\begin{equation}\label{wkr:as:eq:moving}
p=w^3s+B,\qquad q=kw^2s+D,\qquad r=s,
\end{equation}
where
\[
B=b-w^3e,\qquad D=d-kw^2e.
\]
The substitution $(t,v)\mapsto(t,s)$ has Jacobian $H^3$. Accordingly the
old determinant equals $H^3$ times the weighted determinant computed in
$(t,s)$. For arbitrary $A,B,C,D$, the latter is
\begin{align}
&\J(As+B,Cs+D,s)\notag\\
&\quad=DB'-2BD'
+s\bigl(DA'+2CB'-3AD'-2BC'\bigr)
+s^2\bigl(2CA'-3AC'\bigr).\label{wkr:as:eq:movingcoeff}
\end{align}
This identity is just a determinant expansion; the change of variables is
being used only rationally. \textup{[write]} It is Part~I's
\eqref{wkr:eq:threeidentities}, whose coefficients are the left sides of
\eqref{wkr:eq:E0}--\eqref{wkr:eq:E2}; there $A,\ldots,D$ are
polynomials, here rational functions.

Put $A=w^3$, $C=kw^2$, and then write
\[
B=w^2T,\qquad D=wV.
\]
The $s^2$ coefficient vanishes identically. The $s$ coefficient becomes
\[
w^4(2kT'-3V').
\]
It follows that
\begin{equation}\label{wkr:as:eq:VT}
V=\alpha T+\delta,\qquad \alpha=\frac{2k}{3}\ne0,\qquad \delta\in K.
\end{equation}
The constant coefficient satisfies
\[
DB'-2BD'=w^3(VT'-2TV')=w^3(\delta-\alpha T)T'.
\]
Multiplying by $H^3$ and using $\J=-1$ gives
\begin{equation}\label{wkr:as:eq:TU}
-1=U^3(\delta-\alpha T)T'.
\end{equation}
Therefore
\[
L:=\delta-\alpha T
\quad\hbox{satisfies}\quad
LL'=\frac{\alpha}{U^3}.
\]

\subsection{Regularity collapses every higher-degree pole pattern}

Lemma~\ref{wkr:as:lem:pole} applies. Its additional regularity condition
also holds:
\[
D=d-kw^2e=w(2\delta-L)
\]
is regular at any root of $U$, because $d,e$ are polynomials and $H$ is
coprime to $U$. Thus $U$ is linear and $L$ is a nonzero scalar multiple
of $1/U$. Rewriting the constants, we obtain
\begin{equation}\label{wkr:as:eq:eta-theta}
T=\eta+\frac{\theta}{U},\qquad
\delta=\alpha\eta,\qquad
\alpha\theta^2U'=-1.
\end{equation}
In particular $U'\ne0$ and $\theta\ne0$.

This is the all-degree step: the initial numerator polynomials could have
arbitrarily large degree, but the pole equation and regularity force a
single simple pole and a single linear common factor.

\subsection{Polynomiality at the origin forces the remaining slope constant}

Recovering the original constant terms from
\eqref{wkr:as:eq:moving}--\eqref{wkr:as:eq:eta-theta} gives
\begin{align}
b&=\frac{U^3e+H(\eta U^2+\theta U)}{H^3},\label{wkr:as:eq:breconstruct}\\
d&=\frac{kU^2e+(k/3)H(4\eta U+2\theta)}{H^2}.\label{wkr:as:eq:dreconstruct}
\end{align}
Since $b,d$ are polynomials and $(U,H)=1$, these identities imply
\begin{align*}
H^3&\mid U^2e+H(\eta U+\theta),\\
H^2&\mid U^2e+\frac H3(4\eta U+2\theta).
\end{align*}
Subtracting modulo $H^2$ yields
\begin{equation}\label{wkr:as:eq:Hdivides}
H\mid\eta U-\theta.
\end{equation}

We next show $H(0)\ne0$. Multiplying \eqref{wkr:as:eq:breconstruct} by
$H^3$ and evaluating at $0$, a hypothetical $H(0)=0$ would give
$0=U(0)^3e(0)=1$. Set $h_0=H(0)\ne0$. The conditions $b(0)=d(0)=0$,
$U(0)=e(0)=1$ now give
\[
1+h_0(\eta+\theta)=0,\qquad
1+\frac{h_0}{3}(4\eta+2\theta)=0.
\]
Subtracting three times the first equation from the second after clearing
its denominator shows $\eta=\theta$. Both are nonzero. Therefore
\eqref{wkr:as:eq:Hdivides} says $H\mid U-1$. Since $U$ is linear with
$U(0)=1$, $U-1$ is a nonzero scalar multiple of $t$. But $H(0)\ne0$, so
$H$ is constant. Thus $g=H^3$ is constant.

This proves Theorem~\ref{wkr:as:thm:slope} when $g\ne0$; the case $g=0$
was already constant. The algebraic-closure argument proves vanishing of
all positive-degree coefficients of the original $g$, so it descends to
every characteristic-zero field.\qed

\section{Reconstruction and the exceptional branch}
\label{wkr:as:sec:reconstruction}

\subsection{Zero slope forces the tame form}

Suppose $g=0$. Equation $J_1=0$ becomes
\[
e'ac+e(a'c-ac')=0.
\]
If $c\ne0$, Lemma~\ref{wkr:as:lem:constants} applied to $c/(ea)$ shows
$c=\lambda ea$ for a scalar $\lambda$. The same expression holds for
$c=0$ with $\lambda=0$. The constant coefficient equation becomes
\begin{equation}\label{wkr:as:eq:zero-slope}
a\bigl(\lambda e^2b-ed\bigr)'=-1.
\end{equation}
The two factors are polynomials, so both are constant. Since $a(0)=1$,
$a=1$. Integrating and using $b(0)=d(0)=0$ gives
\[
\lambda e^2b-ed=-t.
\]
Consequently $e\mid t$. The condition $e(0)=1$ forces $e=1$, whence
$c=\lambda$ and $d=t+\lambda b$. Together with $b\in t^2K[t]$, this is
exactly \eqref{wkr:as:eq:tame-normal}.
\textup{[write]} This is the argument of Part~I's
Theorem~\ref{wkr:thm:triangular}, case $\gamma=0$, in the present
notation.

\subsection{Nonzero slope gives the explicit noninjective family}

Let $g=\gamma\ne0$. The replacement
\begin{equation}\label{wkr:as:eq:gamma-normalization}
\widetilde p(t,v)=\gamma p(t,v/\gamma),\quad
\widetilde q(t,v)=q(t,v/\gamma),\quad
\widetilde r(t,v)=r(t,v/\gamma)
\end{equation}
preserves polynomiality, the normalized origin coefficients, and the
Jacobian $-1$. It reduces the third numerator to $e+v$.

We may therefore use the proof above with $H=1$. The two origin equations
now give
\[
\eta=\theta=-\frac12.
\]
Put $\beta=e'(0)$. Differentiating \eqref{wkr:as:eq:breconstruct} at
zero and using $b'(0)=0$ gives
\[
0=\frac32U'+\beta.
\]
Thus $U=1-2\beta t/3$ and $\beta\ne0$, since $U'\ne0$. The last identity
in \eqref{wkr:as:eq:eta-theta} gives $kU'=-6$, hence $k=9/\beta$. We
recover
\[
\widetilde p=U^3(e+v)-\frac{U^2+U}{2},\qquad
\widetilde q=\frac9\beta\left(U^2(e+v)-\frac{2U+1}{3}\right).
\]
Undoing \eqref{wkr:as:eq:gamma-normalization} yields
\eqref{wkr:as:eq:normalp}--\eqref{wkr:as:eq:normalr}. Every coefficient is
expressed over the original field, so no algebraic-extension parameter
remains.
\textup{[write]} This is a second proof of Part~I's nonconstant branch
(Section~\ref{wkr:sec:proof}, case $\gamma\ne0$), and of
Theorem~\ref{wkr:thm:triangular}(ii) without the source shear: Part~I
bounds degrees by the weighted Wronskian lemmas, the manuscript by the
pole equation.

\subsection{Converse and uniqueness}

For \eqref{wkr:as:eq:tame-normal}, direct substitution gives $\J=-1$.
For \eqref{wkr:as:eq:normalp}--\eqref{wkr:as:eq:normalr}, first set
$\gamma=1$ as above. Then
\[
A=U^3,\quad C=\frac9\beta U^2,\quad
B=-\frac{U^2+U}{2},\quad D=-\frac3\beta(2U+1).
\]
Equation~\eqref{wkr:as:eq:movingcoeff} has zero $s^2$ and $s$
coefficients, and constant coefficient $-1$, using $U'=-2\beta/3$. This
argument is independent of $e$. Restoring $\gamma$ preserves the
determinant.

At $t=0$ the formulas give $b(0)=d(0)=0$. Differentiation gives
$b'(0)=0$, $d'(0)=1$, and $a(0)=1$, while $e(0)=1$ by hypothesis.
Thus the lift is polynomial and normalized. In (T), the parameters are
$B=p(t,0)$ and $\lambda=q_v$; in (N), they are read from $r$ as stated
above. The branches are disjoint because $r_v=0$ in (T) and $r_v\ne0$
in (N). This proves the algebraic assertions of
Theorem~\ref{wkr:as:thm:classification}.
\section{Equivalence, certificates, and exact degree consequences}
\label{wkr:as:sec:consequences}

\subsection{An explicit inverse on the tame branch}

Write $B=t^2H$. The tame lift of \eqref{wkr:as:eq:tame-normal} is
\eqref{wkr:eq:tameF}, and its inverse is \eqref{wkr:eq:tameinverse}.
Substitution in either order gives the identity. Moreover the map is the
composition of the elementary shear $z\mapsto z+y^2H(xy)$, the elementary
shear $y\mapsto y+\lambda xz$ in the updated coordinates, and the
permutation $(x,y,z)\mapsto(z,y,x)$. Hence it is tame, not merely
invertible.
\textup{[write]} The manuscript displays the lift and the inverse again,
with output coordinates $(X,Y,Z)$ where Part~I writes $(U,V,W)$; they are
printed once, in Part~I. Here $H$ is Part~I's $H$ with $B=t^2H$, not the
cube root of Section~\ref{wkr:as:sec:rigidity}.

\subsection{One core and one source shear on the noninjective branch}

Let $F_{\beta,\gamma}^{\mathrm{core}}$ denote the form (N) with
$e=1+\beta t$, and let $F_*=F_{1,1}^{\mathrm{core}}$. Define
\begin{equation}\label{wkr:as:eq:source-shear}
h(t)=\frac{e(t)-1-\beta t}{\gamma t^2},\qquad
T_h(x,y,z)=\bigl(x,y,z+y^2h(xy)\bigr).
\end{equation}
The numerator defining $h$ vanishes to order at least two, so
$h\in K[t]$. The map $T_h$ is an equivariant elementary shear with
inverse $T_{-h}$, and on invariants it replaces $v$ by $v+t^2h(t)$.
Therefore
\begin{equation}\label{wkr:as:eq:core-shear}
F=F_{\beta,\gamma}^{\mathrm{core}}\circ T_h.
\end{equation}
The core itself satisfies
\begin{equation}\label{wkr:as:eq:diagonal-core}
F_{\beta,\gamma}^{\mathrm{core}}
=\diag(\gamma^{-1},\beta^{-1},1)
\circ F_*\circ\diag(1,\beta,\gamma),
\end{equation}
and $F_*$ is the normalization \eqref{wkr:eq:originalnormalization} of
$F_0$. These are exact polynomial identities. They also explain the
relationship with the previous report's parameters $a_{\mathrm I},
b_{\mathrm I}$: there $\beta=a_{\mathrm I}b_{\mathrm I}$ and
$\gamma=a_{\mathrm I}b_{\mathrm I}^2$, so
$a_{\mathrm I}=\beta^2/\gamma$ and $b_{\mathrm I}=\gamma/\beta$. We avoid
using $a,b$ as family parameters here because they already name the first
numerator coefficients in \eqref{wkr:as:eq:abcdef}.
\textup{[write]} \eqref{wkr:as:eq:source-shear} is Part~I's $h$ of
\eqref{wkr:eq:triangular-normal} with $R=e$, \eqref{wkr:as:eq:core-shear}
is Theorem~\ref{wkr:thm:triangular}(ii), and \eqref{wkr:as:eq:diagonal-core}
is \eqref{wkr:eq:diagonal} with $F^{\mathrm{core}}_{\beta,\gamma}=
F_{a_{\mathrm I},b_{\mathrm I}}$.

\begin{corollary}[Two tame source--target equivalence classes]
\label{wkr:as:cor:twoclasses}
Every map in the stated weighted affine-in-$z$ Keller class is tame
source--target equivalent either to the identity or to $F_0$, and the two
classes are distinct.
\end{corollary}
\begin{proof}
Undo target normalization. The tame branch consists of tame
automorphisms, so lies in the identity class.
Equations~\eqref{wkr:as:eq:core-shear}, \eqref{wkr:as:eq:diagonal-core}
and \eqref{wkr:eq:originalnormalization} put the other branch in the class
of $F_0$. Noninjectivity, shown next, is preserved by source and target
automorphisms and separates the two classes.
\end{proof}

\subsection{Ground-field collision certificates}

The core has the collision
\[
F_{\beta,\gamma}^{\mathrm{core}}(1,0,-1/\gamma)
=F_{\beta,\gamma}^{\mathrm{core}}(0,-9/\beta,71/\gamma)
=(-1/\gamma,-9/\beta,0).
\]
Transporting by $T_{-h}$ gives the following certificate for the full
class.

\begin{proposition}\label{wkr:as:prop:collision}
For every normalized map in (N), the distinct points
\begin{align}
X_1&=(1,0,-1/\gamma),\notag\\
X_2&=\left(0,-9/\beta,\frac{71}{\gamma}-\frac{81h(0)}{\beta^2}\right)
\label{wkr:as:eq:collision}
\end{align}
have common image $(-1/\gamma,-9/\beta,0)$.
\end{proposition}
\begin{proof}
At $t=0$, the invariant numerators are $p=v$, $q=9\gamma v/\beta$,
and $r=1+\gamma v$. This gives the image of $X_1$ directly. On $x=0$, the
core is
\[
(0,y,z)\longmapsto
\left(z-\frac{8\beta^2}{9\gamma}y^2,\ y,\ 0\right).
\]
Putting $y=-9/\beta$, $z=71/\gamma$ gives the same image. The shear does
not move $X_1$ and sends $X_2$ to that core point. The first coordinates
of $X_1$ and $X_2$ differ, so they are distinct over every
characteristic-zero field.
\end{proof}
\noindent\textup{[write]} This is Part~I's collision
\eqref{wkr:eq:newcollision} (with $\gamma=a_{\mathrm I}b_{\mathrm I}^2$,
$\ell=9b_{\mathrm I}=9\gamma/\beta$), made explicit after the shear that
Part~I's proof of Theorem~\ref{wkr:thm:triangular} only mentions; the
mechanism is the project's (Section~\ref{wkr:sub:project}). The
proposition was rechecked at the write with symbolic $\beta,\gamma$ and
$e=1+\beta t+e_2t^2+e_3t^3$.

Unnormalizing the target changes the common output by a diagonal scaling,
but not the colliding source points. Thus the classifier supplies rational
collision certificates whenever the input coefficients are rational.

\begin{example}\label{wkr:as:ex:shear}
Take $\beta=\gamma=1$ and $e=1+t+t^2+t^3$. Then $h=1+t$, and the
collision is
\[
F(1,0,-1)=F(0,-9,-10)=(-1,-9,0).
\]
Its coordinate degrees are $(10,9,7)$, even though its underlying core has
degrees $(7,6,4)$. The additional degree comes entirely from the source
shear, not from a new equivalence class.
\end{example}

\subsection{The exact ordinary degree spectrum}

\begin{corollary}\label{wkr:as:cor:degrees}
For every normalized noninjective map in this class, $h=0$ gives
coordinate degrees $(7,6,4)$ and core support sizes $(7,6,3)$. If $h\ne0$
has degree $m$, the coordinate degrees are exactly
\begin{equation}\label{wkr:as:eq:degreeprofile}
(2m+8,\ 2m+7,\ 2m+5).
\end{equation}
Consequently the possible maximum ordinary degrees of noninjective maps
in the full weighted affine-in-$z$ class are exactly
\[
\{7\}\cup\{8,10,12,14,\ldots\}.
\]
\end{corollary}
\begin{proof}
\textup{[write]} By Theorem~\ref{wkr:as:thm:slope} and
\eqref{wkr:as:eq:core-shear}, every such map is one of Part~I's
Theorem~\ref{wkr:thm:triangular}, so this is
Corollary~\ref{wkr:cor:degreespectrum} together with
Corollary~\ref{wkr:cor:support}. The manuscript's proof repeats Part~I's:
expansion of the core formulas gives the degree and support profiles
($\beta\gamma\ne0$ makes every listed core coefficient nonzero); the
coefficients of $z$ in the three coordinates have degrees $6,5,3$;
substituting $z+y^2h(xy)$ adds nonzero leading terms of degrees
$6+(2+2m)$, $5+(2+2m)$, $3+(2+2m)$, which no old term can cancel; and
$h=t^m$ realizes every $m\ge0$.
\end{proof}

This is a structural exclusion of ordinary degrees $4,5,6$ for
\emph{noninjective maps in this class}, not for all three-variable Keller
maps. Likewise the sixteen-term assertion concerns the unsheared core;
we make no unsupported claim that arbitrary shears preserve total support
size or that sixteen terms are globally minimal.

\section{Function-field consequences and a cubic model}
\label{wkr:as:sec:cubic}

The constant-slope core and its cubic geometry were already part of the
post-announcement discussion reflected in \cite{shaska,gao}. The enlarged
classification now transfers that geometry to every noninjective map of
our class. For completeness we give a direct calculation for $F_0$.

\begin{proposition}\label{wkr:as:prop:cubic}
Let $(P,Q,R)=F_0(x,y,z)$. Then
\[
[K(x,y,z):K(P,Q,R)]=3,
\]
and the Galois group of the normal closure is $S_3$. The same statements
hold after extending $K$ to its algebraic closure, and for every map in
branch \textup{(N)}. In branch \textup{(T)} the degree is $1$.
\end{proposition}
\begin{proof}
The nonzero Jacobian implies algebraic independence of $P,Q,R$ in
characteristic zero. To see this directly, choose a nonzero polynomial
relation of least positive degree if one exists. Differentiating it and
inverting the Jacobian over $K(x,y,z)$ shows that every partial
derivative of the relation is another relation. Minimality makes all
those partial derivatives zero; characteristic zero then makes the
relation constant, a contradiction.

Set $\tau=y+1/x$. Expansion of \eqref{wkr:eq:original} gives
\begin{equation}\label{wkr:as:eq:cubic}
f(\tau)=0,\qquad f(T)=RT^3-2T^2+QT-2P,
\end{equation}
and
\begin{equation}\label{wkr:as:eq:cubic-inverse}
\frac1x=\frac{Q-4\tau+3R\tau^2}{2}=\frac{f'(\tau)}2,
\quad y=\tau-\frac1x,
\quad z=\frac{2x-3x^2y-R}{x^3}.
\end{equation}
Thus $K(x,y,z)=K(P,Q,R,\tau)$.

As a polynomial in $K[P,Q,R,T]$, $f$ is irreducible: it is linear in $P$
with invertible coefficient $-2$, and its quotient ring is the polynomial
ring $K[Q,R,T]$. Viewed as a polynomial in $T$ over $K[P,Q,R]$, it is
primitive because its $T^2$ coefficient is $-2$. Gauss's lemma therefore
gives irreducibility in $K(P,Q,R)[T]$. The extension degree is three.

Its discriminant is
\begin{equation}\label{wkr:as:eq:discriminant}
\Delta=4\bigl(Q^2-RQ^3-16P-27R^2P^2+18PQR\bigr).
\end{equation}
It is not a square in $K(P,Q,R)$. Indeed, a polynomial that is a rational
square over a unique factorization domain is a polynomial square: a
coprime denominator would have to divide the numerator. But specializing
such a square identity at $R=0$ would make $4(Q^2-16P)$ a square in
$K[P,Q]$, impossible because its degree in $P$ is one.

An irreducible separable cubic has a transitive subgroup of $S_3$ as its
normal-closure group. Its discriminant is a square exactly when the group
lies in $A_3$: the product of the three pairwise root differences changes
by the sign of a permutation. The nonsquare discriminant therefore forces
$S_3$. The argument is unchanged over the algebraic closure of $K$.
Source and target polynomial automorphisms preserve these function-field
invariants, so \eqref{wkr:as:eq:core-shear},
\eqref{wkr:as:eq:diagonal-core} and \eqref{wkr:eq:originalnormalization}
transfer them to (N). The explicit inverse gives degree one in (T).
\end{proof}

This statement concerns generic function-field degree. It does not assert
that every special fiber has three points, that the map is finite, or
that its restriction over an unspecified open set is already a covering.

\noindent\textup{[write]} The cubic is not new in ProveIt either. The
sibling report \rpath{jacobian-conjecture/arithmetic-local-global-fibers}
uses the same variable $\tau$ (its $t=y+1/x$), the same cubic
$g_{A,B,C}(T)=CT^3-2T^2+BT-2A$ with $(A,B,C)=(P,Q,R)$, the same
derivative identity $f'(\tau)=2/x$ and the same discriminant (its
Section~2, equation (2.9), equal to \eqref{wkr:as:eq:discriminant} after
expansion);
it credits the cubic inverse model to an archived note of
\texttt{royvanrijn/jacobian-research} and the complex fiber law to Gao's
Theorems~3.3--3.4, and Gao's Theorem~3.3 gives geometric degree three for
this map (Section~\ref{wkr:sub:gao}). The sibling report does not state
that the generic Galois group is $S_3$. The cubic relation, the discriminant identity and the
degree profile of Example~\ref{wkr:as:ex:shear} were rechecked at the
write.

\section{Reduced bases and nilpotent slope directions}
\label{wkr:as:sec:deformations}

Field-valued classification and scheme-theoretic classification are not
interchangeable. This section determines an explicit part of the
difference, rather than assuming that the field theorem automatically
eliminates all coefficient directions.

\subsection{Constancy over every reduced base}

\begin{theorem}[Reduced-base rigidity]\label{wkr:as:thm:reduced}
Let $\Lambda$ be a reduced commutative $\Q$-algebra with identity.
Suppose $a,b,c,d,e,g\in\Lambda[t]$ satisfy
\[
b(0)=b'(0)=d(0)=0,\qquad a(0)=d'(0)=e(0)=1,
\qquad \J=-1.
\]
Then $g\in\Lambda$: every positive-degree coefficient of $g$ vanishes.
\end{theorem}
\begin{proof}
For each prime ideal $\mathfrak p$ of $\Lambda$, reduce the identities
modulo $\mathfrak p$ and pass to $\Frac(\Lambda/\mathfrak p)$, a field of
characteristic zero. Theorem~\ref{wkr:as:thm:slope} makes every
positive-degree coefficient of $g$ zero there, hence zero in the domain
$\Lambda/\mathfrak p$. Thus each such coefficient belongs to every prime
ideal of $\Lambda$, that is, to its nilradical. Reducedness makes it zero.
\end{proof}

We assert only slope constancy over reduced bases, not a global
two-branch normal-form theorem for arbitrary rings. For example, a
constant slope can be a zero divisor on a disconnected base, and division
by it need not be allowed.
\textup{[write]} Part~I's Corollary~\ref{wkr:cor:reduced} is the same
argument for $r=1+t+v$ fixed; here $r$ is free.

\subsection{Explicit dual-number solutions in every slope degree}

Let $\eps^2=0$. For $G\in K[t]$, define
\begin{equation}\label{wkr:as:eq:CG}
C_G(t)=-tG(t)-2\int_0^tG(s)\,ds.
\end{equation}
The integral here is formal polynomial antiderivation in characteristic
zero; in particular $C_G(0)=0$ and
\begin{equation}\label{wkr:as:eq:CGprime}
C_G'=-tG'-3G.
\end{equation}

\begin{theorem}[Nonconstant infinitesimal slopes]\label{wkr:as:thm:dual}
For every $G\in K[t]$, the triple over $K[\eps]/(\eps^2)$
\begin{equation}\label{wkr:as:eq:dual-example}
p=v,\qquad q=t+\eps C_G(t)v,\qquad r=1+\eps G(t)v
\end{equation}
is polynomial, normalized, and satisfies $\J=-1$. If $G$ is nonconstant,
its third invariant slope is a nonconstant polynomial over this
nonreduced coefficient ring.
\end{theorem}
\begin{proof}
Polynomiality and the origin conditions are immediate. Over the
untruncated polynomial ring $K[\eps,t,v]$, formulas
\eqref{wkr:as:eq:J0}--\eqref{wkr:as:eq:J2} give the exact identity
\begin{equation}\label{wkr:as:eq:exact-dual}
\J=-1+\eps^2v^2\mathcal O(G),\qquad
\mathcal O(G)=C_GG'-3GC_G'.
\end{equation}
The first-order term vanishes by \eqref{wkr:as:eq:CGprime}. Modulo
$\eps^2$, the displayed determinant is $-1$.
\end{proof}

For the monomial $G=t^m$ with $m\ge0$, these formulas read
\begin{align}
C_G&=-\frac{m+3}{m+1}t^{m+1},\label{wkr:as:eq:monomialC}\\
\mathcal O(G)&=\frac{(m+3)(2m+3)}{m+1}t^{2m}.\label{wkr:as:eq:monomialO}
\end{align}
When $m\ge1$, the third slope is genuinely nonconstant. In source
coordinates the dual-number map is
\begin{equation}\label{wkr:as:eq:dual-lift}
\left(z,\ y+\eps xz C_G(xy),\ x+\eps x^3zG(xy)\right).
\end{equation}
For $G=t^m$, its maximum ordinary degree is $2m+4$.

\subsection{Every nonzero slope direction is obstructed at second order}

One might try to repair \eqref{wkr:as:eq:exact-dual} by adding
second-order terms. The next theorem rules out \emph{all} such repairs
inside the structural class, not merely the obvious ones.

\begin{theorem}[Universal second-order slope obstruction]
\label{wkr:as:thm:obstruction}
Start from the normalized tame point $(p_0,q_0,r_0)=(v,t,1)$, whose lift
is $(z,y,x)$. Consider any normalized first-order deformation, affine in
$v$,
\begin{align}
p&=v+\eps\bigl(A(t)v+B(t)\bigr),\notag\\
q&=t+\eps\bigl(C(t)v+D(t)\bigr),\label{wkr:as:eq:generalfirst}\\
r&=1+\eps\bigl(E(t)+G(t)v\bigr)\notag
\end{align}
that satisfies $\J=-1$ modulo $\eps^2$. If $G\ne0$, this deformation
has no lift to a Keller solution over $K[\eps]/(\eps^3)$ with all three
numerators still affine in $v$. This remains true with no bound on the
$t$-degrees of the second-order correction polynomials.
\end{theorem}
\begin{proof}
Linearizing \eqref{wkr:as:eq:J0}--\eqref{wkr:as:eq:J2} gives
\begin{align}
[\eps]J_0&=-\bigl(A+D'+E+tE'\bigr),\label{wkr:as:eq:linear0}\\
[\eps]J_1&=-\bigl(C'+tG'+3G\bigr),\label{wkr:as:eq:linear1}\\
[\eps]J_2&=0.\notag
\end{align}
Thus every first-order solution satisfies
\begin{equation}\label{wkr:as:eq:Cconstraint}
C'=-tG'-3G.
\end{equation}
If $G$ has degree $m\ge0$ and leading coefficient $g_m\ne0$, then $C$
has degree $m+1$ and leading coefficient
\[
-\frac{m+3}{m+1}g_m.
\]
The integration constant in $C$ is arbitrary and does not affect this
statement.

Now add arbitrary affine-in-$v$ second-order terms. In the notation of
\eqref{wkr:as:eq:abcdef}, write
\[
a=1+\eps A+\eps^2A_2,\qquad
c=\eps C+\eps^2C_2,\qquad
g=\eps G+\eps^2G_2,
\]
with equally arbitrary $b,d,e$ corrections. The coefficient $J_2$ depends
only on $a,c,g$. Its formula \eqref{wkr:as:eq:J2} gives
\begin{equation}\label{wkr:as:eq:obstruction-universal}
[\eps^2]J_2=CG'-3GC'.
\end{equation}
In particular no second-order correction occurs in this expression, and
neither do $A,B,D,E$.

The leading coefficient of \eqref{wkr:as:eq:obstruction-universal}, in
degree $2m$, is
\begin{equation}\label{wkr:as:eq:leading-obstruction}
\frac{(m+3)(2m+3)}{m+1}g_m^2.
\end{equation}
For $m=0$, the same expression follows directly from $G'=0$ and
$C'=-3G$. Characteristic zero makes \eqref{wkr:as:eq:leading-obstruction}
nonzero. Therefore $J_2$ cannot vanish modulo $\eps^3$, as the Keller
identity would require. This contradiction proves the theorem.
\end{proof}
\noindent\textup{[write]} In this theorem $A,\ldots,E,G$ are first-order
variations and $g_m$ is the leading coefficient of $G$; the coefficients
$g_i$ of Corollary~\ref{wkr:as:cor:schemes} are those of the universal
slope $g$.

The obstruction includes \emph{constant} nonzero $G$, not just the
nonconstant slopes that contradict reduced-base rigidity. It concerns
first-order directions at the particular tame point $(z,y,x)$; it does
not say that the noninjective constant-slope branch is absent from the
coefficient space. A branch can meet a point with its slope vanishing to
higher order in a parameter.

\subsection{The obstruction is sharp at the linearity boundary}

The failure to lift is caused by the affine-in-$v$ restriction, not by an
intrinsic failure of the infinitesimal volume-preserving deformation.
There is an explicit minimal-degree repair.

\begin{proposition}[A quadratic repair]\label{wkr:as:prop:quadratic-repair}
For the canonical direction \eqref{wkr:as:eq:dual-example}, set
\[
I_G(t)=\int_0^t\mathcal O(G)(s)\,ds.
\]
Then, over $K[\eps]/(\eps^3)$,
\begin{equation}\label{wkr:as:eq:quadratic-repair}
p=v,\qquad q=t+\eps C_Gv+\eps^2 I_Gv^2,\qquad
r=1+\eps Gv
\end{equation}
is a normalized polynomial Keller map. For $G\ne0$, maximum $v$-degree
two is the smallest possible for a second-order repair of this direction.
\end{proposition}
\begin{proof}
At $(p,q,r)=(v,t,1)$, the linearized weighted determinant for arbitrary
polynomial variations $(P_1,Q_1,R_1)$ is
\begin{equation}\label{wkr:as:eq:unrestricted-linearization}
-\partial_vP_1-\partial_tQ_1-R_1-t\partial_tR_1-2v\partial_vR_1.
\end{equation}
The second-order addition in \eqref{wkr:as:eq:quadratic-repair} therefore
contributes $-I_G'v^2=-\mathcal O(G)v^2$, canceling
\eqref{wkr:as:eq:exact-dual} modulo $\eps^3$. The extra monomial powers
lift polynomially and do not affect the normalized origin conditions.
Theorem~\ref{wkr:as:thm:obstruction} rules out every degree-one repair
when $G\ne0$, whereas the displayed one has degree two.
\end{proof}
\noindent\textup{[write]} Theorem~\ref{wkr:as:thm:dual} with
\eqref{wkr:as:eq:monomialC}--\eqref{wkr:as:eq:monomialO} and this
repair were rechecked at the write for $G=t^m$, $0\le m\le4$.

\begin{proposition}[Formal integration after removing the degree cap]
\label{wkr:as:prop:formalflow}
Every canonical direction \eqref{wkr:as:eq:dual-example} extends to a
weighted formal Keller family over $K[x,y,z][[\eps]]$. No uniform bound
on the $z$-degree of its coefficients is asserted.
\end{proposition}
\begin{proof}
Consider the polynomial derivation
\[
\mathcal D=x^3zG(xy)\frac{\partial}{\partial x}
          +xzC_G(xy)\frac{\partial}{\partial y},\qquad \mathcal D(z)=0.
\]
Its divergence is $x^2z(3G+tG'+C_G')=0$. It preserves weights, and its
coefficients have no linear terms. The formal exponential
$\Phi_\eps=\exp(\eps\mathcal D)$ is a well-defined automorphism of
$K[x,y,z][[\eps]]$, with inverse $\Phi_{-\eps}$: every coefficient in
$\eps$ is a polynomial, and the exponential identities hold
coefficientwise. The formal flow preserves the Jacobian determinant. More
explicitly, differentiating its Jacobian matrix in $\eps$ and applying
the determinant product rule gives
\[
\frac{\partial}{\partial\eps}\det D\Phi_\eps
  = (\operatorname{div}\mathcal D)\circ\Phi_\eps\,
    \det D\Phi_\eps=0.
\]
The initial determinant is one. Thus
$F_\eps=(z,y,x)\circ\Phi_\eps$ has determinant $-1$, preserves the
required weights, and retains the normalized linear part. Its first-order
term is exactly \eqref{wkr:as:eq:dual-lift}. This is the standard
formal-flow construction applied to the present divergence-free
direction, not a claim of a polynomial family in $\eps$ or analytic
convergence.
\end{proof}

\subsection{Nonzero nilpotents in finite coefficient schemes}

For $\Dcap\ge2$, let $\Lambda_{\Dcap}$ be the universal $\Q$-algebra of
coefficients of six polynomials $a,b,c,d,e,g$ of $t$-degree at most
$\Dcap$, modulo polynomiality, normalization, and the coefficient
equations $\J=-1$. This is a finite presentation: only finitely many
coefficients occur in \eqref{wkr:as:eq:J0}--\eqref{wkr:as:eq:J2}. Write
\[
g(t)=g_0+g_1t+\cdots+g_{\Dcap}t^{\Dcap}
\]
for its universal third slope, and let $\mathfrak m$ be the rational
point corresponding to $(p,q,r)=(v,t,1)$.

\begin{corollary}[Arbitrarily many nilpotent slope directions]
\label{wkr:as:cor:schemes}
For every $\Dcap\ge2$:
\begin{enumerate}[label=\textup{(\roman*)}]
\item Every $g_i$ with $i\ge1$ is nilpotent in $\Lambda_{\Dcap}$.
\item For $1\le i\le \Dcap-1$, $g_i$ is nonzero both in $\Lambda_{\Dcap}$
and in the local ring $(\Lambda_{\Dcap})_{\mathfrak m}$.
\item The tangent space at $\mathfrak m$ contains a
$(\Dcap-1)$-dimensional subspace of nonconstant-slope directions, every
nonzero member of which is obstructed from lifting to order two inside
the affine-in-$v$ class, even if the degree cap is subsequently removed.
\end{enumerate}
\end{corollary}
\begin{proof}
Apply Theorem~\ref{wkr:as:thm:reduced} to $\Lambda_{\Dcap}$ modulo its
nilradical to prove (i). For (ii), the dual-number example with $G=t^i$
uses $C_G$ of degree $i+1$, so lies inside the degree cap when
$i\le \Dcap-1$. Its evaluation homomorphism
$\Lambda_{\Dcap}\to\Q[\eps]/(\eps^2)$ sends $g_i$ to $\eps\ne0$. Since
the residue point is $\mathfrak m$, every element outside $\mathfrak m$
maps to a unit; the homomorphism extends to the local ring. Thus $g_i$ is
also nonzero there.

The first derivatives of these evaluation maps are tangent vectors whose
values on the coordinates $g_1,\ldots,g_{\Dcap-1}$ form the identity
matrix. They are therefore linearly independent. A nonzero linear
combination has a nonzero polynomial $G$ with no constant term, and
Theorem~\ref{wkr:as:thm:obstruction} applies. This proves (iii).
\end{proof}

Thus a field-valued classification does not describe the full coefficient
scheme: nonconstant slopes are invisible at all field-valued points but
survive as genuine nilpotents. We have not determined the nilpotence
indices, all relations among these nilpotents, or the full local algebra.
Those are separate problems rather than consequences of the present
calculation.
\textup{[write]} These are statements at the \emph{tame} point
$(v,t,1)$ with $r$ free. Part~I's Theorem~\ref{wkr:thm:double} fixes
$r=1+t+v$ at the noninjective point; the two coefficient problems are
different, and neither result bears on the other's scheme structure.

\section{An exact recognition algorithm and its verification boundary}
\label{wkr:as:sec:algorithm}

\subsection{Recognition without a determinant search}

Assume the input consists of six densely encoded univariate polynomials
of maximum degree $\Dcap$ over a characteristic-zero field with exact
field operations. The following algorithm decides the Keller condition
inside the stated structural class.

\begin{enumerate}[label=\arabic*.]
\item Check $b(0)=b'(0)=d(0)=0$. A failure means the proposed lift is not
a polynomial map in the class.
\item Compute $a(0)$, $d'(0)$, and $e(0)$. A zero factor implies that the
Jacobian vanishes at the origin, so the map is not Keller. Otherwise
normalize the three outputs by these factors.
\item Check whether $g$ is constant. If not, reject by
Theorem~\ref{wkr:as:thm:slope}.
\item If $g=0$, test $r=1$, $p=v+B$, and $q=t+\lambda(v+B)$ with
$B=p(t,0)$ and $\lambda=q_v(0)$. On success, return the inverse
\eqref{wkr:eq:tameinverse} and undo output normalization.
\item If $g=\gamma\ne0$, compute $\beta=e'(0)$. Reject if $\beta=0$;
otherwise compare the input with
\eqref{wkr:as:eq:normalp}--\eqref{wkr:as:eq:normalr}. On success return
$\beta,\gamma,e,h$ and the collision \eqref{wkr:as:eq:collision}.
\end{enumerate}

Theorem~\ref{wkr:as:thm:classification} proves correctness in both
directions. No polynomial factorization, extraction of roots, algebraic
closure, or unbounded search is required by the algorithm. Those tools
are used in the proof of the theorem, not by an implementation of the
normal-form test.
\textup{[write]} This extends Part~I's procedure of
Section~\ref{wkr:sec:verification}, whose scope stopped at a constant
coefficient of $v$ in $r$.

\begin{proposition}[Arithmetic-operation complexity]
\label{wkr:as:prop:complexity}
The normal-form recognition algorithm uses $O(\Dcap)$ field operations
for dense univariate coefficient arrays, in addition to reading the
input.
\end{proposition}
\begin{proof}
All origin tests, scalar normalizations, derivative coefficient reads,
and coefficient comparisons take linear time. The only polynomial
multiplications needed in (N) are by $U^2$ and $U^3$, of fixed degrees
two and three. The shift defining $h$ is a subtraction followed by scalar
and monomial division. Each operation is linear in the input degree.
\end{proof}

This is an arithmetic model bound, not a bit-complexity bound over $\Q$.
Nor is it a performance guarantee for a general-purpose symbolic
expression engine. The companion SymPy program implements the same tests
exactly but uses expression expansion and polynomial conversion for
clarity.

\subsection{The executable companion}

The package contains \rpath{code/weighted_keller.py}, restricted to
rational input coefficients. Its public functions provide weighted
Jacobians, polynomial lifting, normalized noninjective forms,
classification, collision points, and infinitesimal obstruction
polynomials. Inputs outside the structural class raise a descriptive
\texttt{ValueError}; a valid polynomial input in the class that is not
Keller returns a negative classification.

For example:
\begin{lstlisting}[language=Python]
from code.weighted_keller import t, normal_form, classify
p, q, r = normal_form(1, 1, 1 + t + t**2 + t**3)
result = classify(p, q, r)
assert result.kind == "noninjective"
\end{lstlisting}
The executable implementation accepts rational coefficients rather than
arbitrary abstract characteristic-zero fields. This restriction belongs
to the software, not to the theorem.
\textup{[write]} In ProveIt the module is shipped as
\rpath{code/02-affine-slope-weighted_keller.py}, which is not an
importable module name; the example works in a copy with the delivered
names (Section~\ref{wkr:as:sec:rerun}).

\subsection{Twelve exact-check groups}

The script \rpath{code/verify_results.py} was run with SymPy~1.14.0.
It records its Python and SymPy versions and all results in
\rpath{data/verification.json}, with a compact text summary alongside.
Every group passed:
\begin{center}
\small
\begin{tabular}{@{}p{0.08\textwidth}p{0.84\textwidth}@{}}
\toprule
Group & Exact assertion checked\\
\midrule
1 & The three generic coefficients of the weighted determinant.\\
2 & The rational moving-coordinate and pole-equation reductions.\\
3 & The noninjective normal form with symbolic $\beta,\gamma,e(t)$.\\
4 & $\det DF_0=-2$, the integral collision, and the exact normalization.\\
5 & Rational diagonal equivalence and source-shear examples.\\
6 & Both tame inverse identities for a symbolic function $H$.\\
7 & Eighteen noninjective rational cases, eighteen rejected mutations,
 twelve tame cases, and three exceptional rejections.\\
8 & The core degree/support profiles and shears of degrees $0$ through
 $10$.\\
9 & Generic first-order equations and independence of the second-order
 obstruction from all second-order affine-in-$v$ corrections.\\
10 & Dual-number examples $G=t^m$ for $0\le m\le12$ and three mixed
 polynomials; the quadratic second-order repair.\\
11 & Cubic elimination, the inverse $x$ formula, and the discriminant.\\
12 & Input-domain, polynomiality, and singular-origin validation.\\
\bottomrule
\end{tabular}
\end{center}
The randomized rational examples use the fixed seed $20260929$. Generic
symbolic identities are distinguished from finite examples in the JSON
output. The root-multiplicity, pole, nilradical, and all-degree arguments
are written proofs; finite verification is not substituted for them.
\textup{[write]} The paths in this subsection are the delivered ones; in
ProveIt they are \rpath{code/02-affine-slope-verify_results.py} and
\rpath{data/02-affine-slope-verification.json} (and \texttt{.txt}).
Part~I's own script is the unprefixed \rpath{code/verify_results.py}.
Group~4 rechecks kernel-checked statements
(Section~\ref{wkr:as:sec:formal}).
\section{Further research questions and concrete next steps}
\label{wkr:as:sec:future}

The original nonconstant-slope question is settled by
Theorem~\ref{wkr:as:thm:slope}; it should not remain on the open-problem
list under the same hypotheses. The following extensions remain distinct.

\subsection{Quadratic dependence on the second invariant}

Allow $\deg_v p,\deg_v q\le2$, initially retaining $\deg_v r\le1$.
Must the third slope still be constant? If not, classify the smallest
new families up to equivariant tame equivalence. The present pole
argument uses both the linear moving-coordinate expansion and the
exponent pattern in $J_2$; neither automatically survives an added $v^2$
coefficient. A useful first milestone is an all-degree description of the
highest two coefficient equations, rather than a fixed ordinary-degree
search alone. This is the next linearity extension proposed in Part~I
(Section~\ref{wkr:as:q:quadratic}).

\subsection{How does leaving the affine class remove the obstruction?}

Proposition~\ref{wkr:as:prop:quadratic-repair} settles the second-order
question sharply: a quadratic term suffices and an affine repair is
impossible. Proposition~\ref{wkr:as:prop:formalflow} also gives
unrestricted formal integration. What remains is the bounded-degree
problem: can a given nonzero direction be continued to every order with a
fixed bound, such as two, on the $v$-degree? Which directions lie on an
actual polynomial parameter family rather than only a formal one with
increasing degrees? These questions cannot be answered merely by taking
the formal exponential.

\subsection{The local nilpotent algebra at the tame boundary}

For the finite coefficient rings $\Lambda_{\Dcap}$, determine the
nilpotence indices of $g_1,\ldots,g_{\Dcap}$ and their mutual products.
Corollary~\ref{wkr:as:cor:schemes} shows that many are nonzero nilpotents
and supplies tangent directions; it does not identify a complete
local-ring presentation. A concrete task is to compute a small $\Dcap$
presentation, isolate relations forced by the universal obstruction, and
prove a pattern by explicit ideal-membership certificates. In particular,
does the nilpotence order stay bounded as $\Dcap$ increases?

\subsection{Higher-degree thickening at a noninjective point}
\label{wkr:as:sec:future-thick}

Part~I gives a double-point coefficient scheme under an ordinary
degree-seven cap with $r=1+t+v$ (Theorem~\ref{wkr:thm:double}), and asks
whether that thickening stabilizes as the degree cap increases
(Section~\ref{wkr:as:q:stabilize}). Our theorem about slope directions at
$(v,t,1)$ does not answer this different question at a noninjective
point. Solve the untruncated linearized system there, then its
obstruction equations, while keeping the fixed-third-numerator condition
explicit.

\subsection{Scheme-theoretic gluing of the branches}
\label{wkr:as:sec:future-glue}

Field-valued normal forms and tame equivalence do not determine how the
noninjective branch approaches the tame locus inside coefficient space.
The second-order theorem is a useful constraint: at the simplest tame
point, a nonzero first-order slope cannot occur along an honest
second-order deformation within the affine class. Determine which higher
orders of slope vanishing occur on actual arcs, and how they depend on
the remaining first jets. This should be compared with the
parameter-plane and boundary calculations in Part~I
(Theorem~\ref{wkr:thm:plane}, Section~\ref{wkr:as:q:gluing}), without
equating a reduced parameterization with the full coefficient scheme.

\subsection{Other mixed weights and generalized pole equations}

Replace $(-1,1,2)$ by $(-1,1,m)$, or more general primitive mixed-sign
triples compatible with a nonsingular linear term. The highest
determinant coefficient should give different power relations and
equations resembling $L^rL'=\alpha/U^s$. Which exponent pairs force a
single critical root, and which permit several? A classification should
track polynomiality along the coordinate hyperplane, not just solve the
localized equations. The arithmetic of the pole orders may provide a
systematic dividing line between rigid and flexible weighted classes.
\textup{[write]} Part~I's Section~10.5 asks the analogous question for
its Wronskian exponents; here $r$ and $s$ in $L^rL'$ are exponents.

\subsection{Positive characteristic under degree hypotheses}

Several proof steps fail in characteristic $p$: derivatives of $p$th
powers vanish, multiplicity drops can change, and coefficients such as
$(m+3)(2m+3)/(m+1)$ need not be nonzero. Determine explicit degree ranges
and excluded primes for which versions of constant-slope rigidity and
the obstruction theorem remain valid. Then construct and classify the
first Frobenius-driven failures. Simply reducing the characteristic-zero
formulas does not prove an unrestricted positive-characteristic theorem.

\subsection{A formal verification route}

A compact formalization can be organized around the actual dependencies:
weighted monomial encoding and localization; the three determinant
coefficients; rational constants of differentiation; valuation
factorization; the one-critical-root pole lemma; polynomial divisibility
and origin constraints; and reconstruction. The nilpotent obstruction
then reduces to coefficient extraction and a leading-term argument.
The companion symbolic identities are suitable targets for kernel-checked
normalization certificates, but a proof assistant must still verify the
all-degree pole and factorization steps. No Lean or Rocq file in this
package is presented as already proving them.

\subsection{Recognition, sparsity, and arithmetic optimization}

The coefficient-only test separates structural classification from
coefficient optimization. Over $\Q$, one can now optimize denominator
height, integral coefficient height, or collision height within the
explicit diagonal-and-shear orbit. Such a project should specify its
normalization and cost function: minimizing ordinary degree, number of
expanded monomials, coefficient height, and the size of a formal
certificate are different problems. Extending the software to sparse
arrays and exact algebraic number fields is also useful, but does not
change the field theorem.

\section{Conclusion of Part~II}
\label{wkr:as:sec:conclusion}

The passage from a scalar third slope to a polynomial third slope was a
real missing case in the repository's weighted classification. It is now
excluded without a degree cap. The mechanism is a rational differential
invariant whose poles are rigid enough that global polynomiality forces
the slope constant. This completes the field-valued classification of
all affine-in-$z$ Keller maps with the specified weights and turns it
into a direct recognition procedure with concrete inverse or collision
certificates.

At the same time, the coefficient schemes retain information that no
field-valued normal-form theorem can see. Nonconstant slopes exist over
dual numbers in arbitrarily high degrees, and every nonzero slope
direction at the simplest tame point has a computable, unavoidable
second-order obstruction inside the class. The combination of exact
rigidity and explicit nilpotent failure gives both a completed result and
a sharply defined next research program.

\section{Proof audit and dependency ledger}
\label{wkr:as:app:audit}

\textup{[write]} This section is the manuscript's Appendix~A.

\begin{center}
\small
\begin{tabular}{@{}p{0.28\textwidth}p{0.64\textwidth}@{}}
\toprule
Claim & What the proof uses\\
\midrule
Constant slope & Characteristic zero; all three numerators affine in $v$;
 polynomial lift; unit normalized origin factors; rational factorization
 over an algebraic closure; regularity at roots of $U$; origin
 constraints.\\
Full classification & Constant-slope theorem; the explicit zero-slope
 calculation; reconstruction and determinant expansion.\\
Tame/noninjective split & Written inverse and two explicitly distinct
 ground-field source points; no topological argument.\\
Degree spectrum & Exact core leading terms and a nonzero source-shear
 leading term; no enumeration.\\
Cubic degree and $S_3$ & Direct identities; algebraic independence from
 nonzero Jacobian; Gauss's lemma; discriminant criterion for a cubic.\\
Reduced-base rigidity & Field theorem in residue fraction fields;
 intersection of prime ideals equals the nilradical.\\
Second-order obstruction & Generic determinant coefficient $J_2$;
 first-order slope equation; nonzero leading coefficient.\\
Coefficient-scheme nilpotents & Reduced-base rigidity; dual-number
 evaluation maps; their linearly independent tangent values.\\
\bottomrule
\end{tabular}
\end{center}

\paragraph{Potential gaps explicitly avoided.}
The moving coordinate is used in a rational function field, not declared
a global automorphism. The root argument is performed over an algebraic
closure and descended only after showing coefficient vanishing. The case
$c=0$ is ruled out before division by $c$, and the case $g=0$ is treated
separately. A constant Jacobian is not inferred from finite examples.
The reduced-base theorem is not extended without qualification to
nonreduced rings. The obstruction theorem allows arbitrary second-order
coefficients rather than testing only the displayed family. Ordinary
polynomial degree is kept separate from geometric degree.

\paragraph{What has not been certified.}
The written proofs have not been checked by a proof-assistant kernel or
an independent referee. The exact computation checks do not establish
literature priority, a global minimum degree, a classification beyond the
linearity and weight hypotheses, or a presentation of the entire
coefficient scheme. The source comparison is restricted to the materials
listed in the references and in \rpath{SOURCES.md}
(\textup{[write]} shipped as \rpath{02-affine-slope-SOURCES.md}).

\section{Reproducibility of Part~II}
\label{wkr:as:sec:repro}

\textup{[write]} The first paragraph is the manuscript's Appendix~B,
with the delivered paths.

The package contains the complete \LaTeX{} source, the compiled PDF,
\rpath{code/weighted_keller.py}, \rpath{code/verify_results.py},
\rpath{requirements.txt}, and the recorded check output. From the
package root:
\begin{lstlisting}[language=bash]
python -m pip install -r requirements.txt
python code/verify_results.py
latexmk -pdf -interaction=nonstopmode -halt-on-error article.tex
\end{lstlisting}
The article uses standard \LaTeX{} packages and embedded Computer Modern
family fonts through Latin Modern. No external bibliography processor or
network access is required to rebuild the PDF once the packages are
installed. The verification scripts perform exact rational and symbolic
arithmetic; they do not download or modify the repository.

\subsection{Shipped names}\label{wkr:as:sec:shipped}
\textup{[write]} In ProveIt the package's companion files carry the
prefix \texttt{02-affine-slope-}, the next number in this report's local
sequence (Part~I's files are unprefixed):
\begin{center}
\small
\begin{tabular}{@{}p{0.40\textwidth}p{0.52\textwidth}@{}}
\toprule
Delivered & Shipped\\
\midrule
\rpath{code/verify_results.py} & \rpath{code/02-affine-slope-verify_results.py}\\
\rpath{code/weighted_keller.py} & \rpath{code/02-affine-slope-weighted_keller.py}\\
\rpath{code/__init__.py} & \rpath{code/02-affine-slope-__init__.py}\\
\rpath{data/verification.json}, \rpath{.txt} &
\rpath{data/02-affine-slope-verification.json}, \rpath{.txt}\\
\rpath{requirements.txt} & \rpath{data/02-affine-slope-requirements.txt}\\
\rpath{SOURCES.md}, \rpath{STATUS.md} &
\rpath{02-affine-slope-SOURCES.md}, \rpath{02-affine-slope-STATUS.md}\\
\bottomrule
\end{tabular}
\end{center}
They are byte-identical to the delivery. The manuscript, its README and
its PDF are not shipped: this Part is the manuscript, and the report's
\rpath{article.pdf} is a build of this text.

\subsection{Rerunning the checks}\label{wkr:as:sec:rerun}
\textup{[write]} Do not run the prefixed script in place. It imports its
helper by the delivered name (\texttt{from weighted\_keller import}),
which does not exist under the shipped names, so it stops with an
\texttt{ImportError}; and if that import were satisfied it would write
\rpath{data/verification.json} and \rpath{data/verification.txt} next to
the \rpath{code/} directory, which in this report are \emph{Part~I's}
records. The package's \rpath{__init__.py} also makes \rpath{code/} a
package that shadows Python's standard module \texttt{code}. Rerun it on
a copy with the delivered names instead:
\begin{lstlisting}[language=bash]
mkdir -p <scratch>/pkg/code
cp code/02-affine-slope-verify_results.py <scratch>/pkg/code/verify_results.py
cp code/02-affine-slope-weighted_keller.py <scratch>/pkg/code/weighted_keller.py
cp code/02-affine-slope-__init__.py <scratch>/pkg/code/__init__.py
cp data/02-affine-slope-requirements.txt <scratch>/pkg/requirements.txt
cd <scratch>/pkg && python code/verify_results.py
\end{lstlisting}
It writes \rpath{data/verification.json} and \rpath{data/verification.txt}
inside the copy; compare them with the shipped
\rpath{data/02-affine-slope-verification.json} and \rpath{.txt}. At the
write this was done from the shipped files with SymPy~1.14.0 (under
\texttt{uv}, Python~3.13.5) in about 11~seconds: all twelve groups passed
and both records equal the shipped ones apart from line endings. The JSON
embeds the Python version, so a run under another version differs in
that field.

%% End of Part II. The appendices below belong to Part I.
\clearpage
\fancyhead[L]{\small\color{inkblue}All-degree rigidity of a weighted Keller class}
\phantomsection\addcontentsline{toc}{part}{Appendices to Part~I}
\appendix
\section{An explicit presentation of the coefficient ideal}
\label{wkr:app:presentation}

For readers who prefer generators to the vector formula
\eqref{wkr:eq:pointplus}, Theorem~\ref{wkr:thm:double} says that $I$ is exactly
the ideal generated by
\begin{equation}\label{wkr:eq:idealexplicit}
\begin{aligned}
&q_{40},\qquad (q_{21}-4)^2,\\
&144p_{20}+7q_{21}+100,\\
&8p_{11}+q_{21}+12,\\
&216p_{30}-27q_{21}-116,\\
&6p_{21}-q_{21}-4,\\
&54p_{40}+3q_{21}+4,\\
&54p_{31}+3q_{21}+4,\\
&8q_{01}-9q_{21}-36,\\
&8q_{20}+11q_{21}+20,\\
&4q_{11}+9q_{21}+12,\\
&q_{30}-q_{21}.
\end{aligned}
\end{equation}
These equations apply over every $\Q$-algebra, not just over fields.
Over a field the square forces $q_{21}=4$, and all remaining coefficients
are the point \eqref{wkr:eq:cstar}. Over the dual numbers, $q_{21}=4+\eps$
produces exactly the tangent vector \eqref{wkr:eq:u}.

As a useful independent check of the equation reconstruction, the
coefficients corresponding to the last three monomials $t,v^2,v$ in
\eqref{wkr:eq:monomialorder} are, respectively,
\begin{align*}
f_{16}&=-p_{11}+2p_{20}q_{01}-2q_{20}-2,\\
f_{17}&=2p_{11}q_{01}-3q_{11},\\
f_{18}&=p_{11}q_{01}+q_{01}-q_{11}-3.
\end{align*}
For the first two monomials they are
\[
f_1=-5p_{31}q_{40}-4p_{40}q_{40},\qquad
f_2=-9p_{31}q_{40}.
\]
The remaining equations are recovered by the same determinant expansion;
the complete list is included in the JSON certificate.

\section{Small exact certificates for the deformation calculation}
\label{wkr:app:certificates}

The primitive integral tangent vector is $n=144u$. To keep the second-order
residual integral, define
\[
\widehat H_i=[e^2]f_i(\cstar+en)=144^2 H_i.
\]
In the ordering \eqref{wkr:eq:monomialorder}, the nonzero data and the
cokernel functional are as follows. Zeros are included to make the ordering
unambiguous.
\begin{center}
\begin{tabular}{@{}rcrr@{}}
\toprule
Row & Monomial & $\widehat H_i$ & $L_i$\\
\midrule
1 & $t^7$ & 0 & $-205$\\
2 & $t^6v$ & 0 & 0\\
3 & $t^6$ & 0 & $-46$\\
4 & $t^5v$ & 0 & 0\\
5 & $t^5$ & 1008 & $-40$\\
6 & $t^4v$ & 1872 & 24\\
7 & $t^4$ & $-4140$ & $-8$\\
8 & $t^3v^2$ & 864 & 0\\
9 & $t^3v$ & $-9504$ & 4\\
10 & $t^3$ & 2772 & 0\\
11 & $t^2v^2$ & $-5184$ & 0\\
12 & $t^2v$ & 12636 & 0\\
13 & $t^2$ & 1620 & 0\\
14 & $tv^2$ & 9720 & 0\\
15 & $tv$ & $-2592$ & 0\\
16 & $t$ & $-2268$ & 0\\
17 & $v^2$ & $-5832$ & 0\\
18 & $v$ & $-2916$ & 0\\
\bottomrule
\end{tabular}
\end{center}
The obstruction is visible in four integer products:
\[
(-40)(1008)+24(1872)+(-8)(-4140)+4(-9504)=-288.
\]
Dividing by $144^2$ gives $LH=-1/72$.

For the nonzero minor in \eqref{wkr:eq:minorvalue}, the zero-based row indices
in the JSON matrix are
\[
(0,2,4,5,6,9,11,12,13,15,17),
\]
and the zero-based column indices are
\[
(0,1,2,3,4,5,6,7,8,9,11).
\]
Thus an independent checker only needs the reconstructed derivative
matrix, one $11\times11$ determinant, and the two matrix-vector products
$Mu$ and $LM$. The companion JSON stores the matrix itself as rational
strings for convenient cross-checking.

\section{Reproduction and source ledger}
\label{wkr:app:reproduction}

The supplied archive contains the following files:
\begin{center}
\begin{tabular}{@{}p{0.36\textwidth}p{0.56\textwidth}@{}}
\toprule
File & Purpose\\
\midrule
\texttt{article.tex} & Self-contained LaTeX source, including bibliography.\\
\texttt{article.pdf} & Compiled article.\\
\texttt{verify\_results.py} & Reconstructs and checks all finite identities.\\
\texttt{verification.json} & Exact equations, matrix, and certificate data.\\
\texttt{verification.txt} & Output of the successful verification run.\\
\texttt{README.md} & Build commands, scope, and trust boundary.\\
\texttt{requirements.txt} & Reproducible SymPy dependency.\\
\bottomrule
\end{tabular}
\end{center}
Build the article with two runs of \texttt{pdflatex article.tex}, or with
\texttt{latexmk -pdf article.tex}. Run the checks with
\begin{quote}\ttfamily
python verify\_results.py -{}-output verification.json
\end{quote}
\wtag{} In ProveIt the script is \rpath{code/verify_results.py}; the
records and \rpath{requirements.txt} are in \rpath{data/}; the delivered
PDF is not shipped (\rpath{article.pdf} is a build of this text); the
delivered README is replaced by the report README. The command above is
the delivery's; its default output \rpath{verification.json} is relative
to the working directory, so in ProveIt pass an output path outside the
report (Appendix~\ref{wkr:app:rerun}).

The verification does not download or execute any repository code. The
repository is a cited mathematical source; its original map and the
relevant formulas are reconstructed in the script. The source ledger is:
\begin{center}
\begin{tabular}{@{}p{0.30\textwidth}p{0.61\textwidth}@{}}
\toprule
Source category & Role in the report\\
\midrule
Repository map & Original map, its known collision, and the finite-search
baseline; not new claims of this article.\\
External research & Context for graded maps and examples outside our class;
not dependencies of the classification proof.\\
Standard algebra & Nullstellensatz and Nakayama's lemma in the passage from
an isolated obstructed tangent to a double point.\\
This article & All-degree classification, family closure, reduced rigidity,
and the full degree-seven coefficient-algebra identification.\\
\bottomrule
\end{tabular}
\end{center}

\section{Provenance of this report}\label{wkr:app:report}

\subsection{One manuscript}\label{wkr:app:onesource}
\wtag{} This report is built from \emph{one} manuscript: manuscript~01 of
batch~36 of ProveIt's incoming reports, the archive
\texttt{ProveIt\_Weighted\_Keller\_Rigidity} (inner directory of the same
name, main file \texttt{article.tex}, a 23-page PDF dated 24~September
2026, PDF metadata author ``ChatGPT; prepared for Vladimir Reshetnikov''),
pinned to ProveIt commit
\nolinkurl{e21766d04c2b8a9b2cdba0cd43e563024b3bd1b9}. The archive arrived in
commit \texttt{e13affd32} and was placed in \texttt{1a1396d4d} as a new
report. Nothing was merged and nothing was selected away: every result,
proof, example, table, question and limitation of the manuscript is
printed. The manuscript contributed everything except the text marked
\wtag.
[Added 29 September 2026, batch~44: this appendix describes Part~I. The
report is now built from \emph{two} manuscripts; Part~II is batch-44
manuscript~04, pinned to \nolinkurl{9b24a3a8d545af9624f6ac455f5b548be62818b6},
and its provenance, shipped files and editorial changes are in
Sections~\ref{wkr:as:sec:source}, \ref{wkr:as:sec:duplicates}
and~\ref{wkr:as:sec:repro}.]

The editorial changes are these.
\begin{enumerate}[label=(\arabic*)]
\item Every label carries the prefix \texttt{wkr:}. The source's 96 labels
are kept, unchanged after the prefix. Six were added:
\texttt{wkr:sub:project}, \texttt{wkr:sub:gao}, \texttt{wkr:app:report},
\texttt{wkr:app:onesource}, \texttt{wkr:app:pinned} and
\texttt{wkr:app:rerun}.
\item Sections~\ref{wkr:sub:project} and~\ref{wkr:sub:gao} are new and
end Section~\ref{wkr:sec:intro}, so no source number moved. They credit
the project's weighted determinant identity and collision mechanism,
state the project's formal status, and correct the reading of Gao's
``degree four''.
\item Notes marked \wtag{} were added on the title page, after the proof
of Proposition~\ref{wkr:prop:jacobian}, after the collision family
\eqref{wkr:eq:collisionfamily}, in Sections~\ref{wkr:sec:verification}
and~\ref{wkr:sec:future}, and in Appendix~\ref{wkr:app:reproduction}; this
appendix is new. The PDF bookmark of the title page was fixed (no
duplicate page anchor), and the command-line option
\texttt{-{}-output} in Appendix~\ref{wkr:app:reproduction} is now typeset
with two hyphens (the delivered PDF showed one dash).
\end{enumerate}
No mathematical statement of the manuscript was changed and no symbol was
renamed; the added notes write $\alpha$ for the project's coefficient
$a=[t^2]p$, because $a$ is a family parameter here. The report's $t=xy$
agrees with the project's research notes but not with the sibling report
on arithmetic fibers, where $t=y+1/x$.

\subsection{Repository statements at the pin}\label{wkr:app:pinned}
\wtag{} The source read the project README, its research notes and
\rpath{verify_equivariant_sparsity.py} at its pin. The project did not
change between the pin and the placement commit. Its statements about the
project are accurate: the map, determinant and collision are the
project's (and formalized, Section~\ref{wkr:sub:project}); the sparsity
script proves that eleven coefficients cannot vanish and identifies the
branch $q_{40}=0$ (script lines 110--127 and 164--167); its certificate
contains $(q_{21}-4)^2$ with the comment that the square ``records a
harmless nonreduced structure'' (lines 129--132), and neither README
identifies the scheme as a double point. The one omission, the credit for
the determinant identity and the collision mechanism, is repaired in
Section~\ref{wkr:sub:project}. Shaska \cite{shaska} and the Stacks Project
references were not checked at the write; Gao \cite{gao} was checked for
Theorems~3.3 and~3.5 (Section~\ref{wkr:sub:gao}).

At the write, the ideal statement of Theorem~\ref{wkr:thm:double} was also
checked in both directions, which the shipped script does not do: with
SymPy~1.14.0 the eighteen coefficient equations and the twelve generators
of \eqref{wkr:eq:idealexplicit} have the same reduced Gr\"obner basis
(grevlex), and each set lies in the ideal of the other. The normalization
\eqref{wkr:eq:originalnormalization} was rechecked the same way.

\subsection{Rerunning the checks}\label{wkr:app:rerun}
\wtag{} The script writes its JSON record to \texttt{-{}-output}, default
\rpath{verification.json} in the current directory. Run it from the report
with an output path outside the repository, so that neither the report
root nor \rpath{data/} is written:
\begin{quote}\ttfamily
python code/verify\_results.py -{}-output <scratch>/verification.json
\end{quote}
It needs Python~3.10 or later and SymPy (\rpath{data/requirements.txt}
pins \texttt{sympy==1.14.0}). At the write it was run so with Python~3.13.5
and SymPy~1.14.0: all thirteen check groups passed, the JSON record equals
\rpath{data/verification.json} apart from line endings, and the printed
output equals \rpath{data/verification.txt} apart from line endings and the
certificate path in its last line. The record embeds the Python version, so
a run under another version differs in that field.

\begin{thebibliography}{9}
\small
\bibitem{repo-main}
Vladimir Reshetnikov, \emph{ProveIt: Algebra/JacobianConjecture/README.md}.
Repository snapshot \texttt{e21766d04c2b8a9b2cdba0cd43e563024b3bd1b9},
inspected September 24, 2026.
\url{https://github.com/VladimirReshetnikov/ProveIt/blob/e21766d04c2b8a9b2cdba0cd43e563024b3bd1b9/Algebra/JacobianConjecture/README.md}.

\bibitem{repo-research}
Vladimir Reshetnikov, \emph{ProveIt: Algebra/JacobianConjecture/Research/README.md}.
Same repository snapshot and inspection date.
\url{https://github.com/VladimirReshetnikov/ProveIt/blob/e21766d04c2b8a9b2cdba0cd43e563024b3bd1b9/Algebra/JacobianConjecture/Research/README.md}.

\bibitem{repo-sparsity}
Vladimir Reshetnikov, \emph{verify\_equivariant\_sparsity.py}.
Same repository snapshot and inspection date.
\url{https://github.com/VladimirReshetnikov/ProveIt/blob/e21766d04c2b8a9b2cdba0cd43e563024b3bd1b9/Algebra/JacobianConjecture/Research/verify_equivariant_sparsity.py}.

\bibitem{repo-9b24}
Vladimir Reshetnikov and contributors,
\emph{ProveIt: Algebra/JacobianConjecture}, repository development and
research notes, inspected at commit
\nolinkurl{9b24a3a8d545af9624f6ac455f5b548be62818b6},
29 September 2026 (cited by Part~II).
\url{https://github.com/VladimirReshetnikov/ProveIt/tree/9b24a3a8d545af9624f6ac455f5b548be62818b6/Algebra/JacobianConjecture}.
The formal status asserted by the repository concerns its own
declarations, not the new theorems of Part~II.

\bibitem{shaska}
T.~Shaska, \emph{Graded Keller maps and the Jacobian Conjecture},
arXiv:2607.20210, 2026. Version 2 dated July 25, 2026;
consulted September 24, 2026.
\url{https://arxiv.org/abs/2607.20210}.
Part~II consulted the HTML of the same version on September 29, 2026,
for graded coordinates, the known counterexample and coefficient-scheme
geometry.

\bibitem{gao}
Shuhong Gao, \emph{Counterexamples to the Jacobian conjecture in dimensions
greater than two}, arXiv:2608.00222, 2026. Version 1 dated July 31, 2026;
consulted September 24, 2026.
\url{https://arxiv.org/abs/2608.00222}.
Part~II consulted the HTML of the same version on September 29, 2026,
for broader constructions and generic fiber degrees; it is not used there
as a classification theorem.

\bibitem{nullstellensatz}
The Stacks Project Authors, \emph{The Stacks Project},
Theorem 10.34.1, Hilbert Nullstellensatz, Tag 00FV.
Consulted September 24, 2026.
\url{https://stacks.math.columbia.edu/tag/00FV}.

\bibitem{nakayama}
The Stacks Project Authors, \emph{The Stacks Project},
Lemma 10.20.1, Nakayama's lemma, Tag 00DV.
Consulted September 24, 2026.
\url{https://stacks.math.columbia.edu/tag/00DV}.
\end{thebibliography}
\end{document}
